\documentclass[11pt]{article}
\usepackage{../common/pbstyle}
%% Common header for TD / TP / project sheets (same layout as the other sheets of the course)
\newcommand{\sheethead}[2]{%
\begin{center}
{\LARGE\bfseries Ensemble Methods}\\[0.2cm]
{\Large #1}\\[0.3cm]
{\large M2 Computer Science -- MALIA}\\[0.4cm]
Guillaume Metzler\\
Institut de Communication (ICOM)\\
Universit\'e de Lyon, Universit\'e Lumi\`ere Lyon 2\\
Laboratoire ERIC UR 3083, Lyon, France\\[0.2cm]
\texttt{guillaume.metzler@univ-lyon2.fr}
\end{center}
\vspace{0.3cm}
\begin{center}\begin{minipage}{0.88\linewidth}\small\textbf{Abstract.} #2\end{minipage}\end{center}
\vspace{0.3cm}}

\renewcommand{\course}{Ensemble Methods -- TD PAC-Bayes}
%% \soltrue is set by the solution wrapper (td_pacbayes_solutions.tex)
\providecommand{\ifsol}{\iffalse}
\usepackage{comment}
\ifsol
\newtcolorbox{solution}{enhanced,breakable,colback=pbteal!4,colframe=pbteal,boxrule=0.4pt,sharp corners,title={Solution},fonttitle=\bfseries\small,colbacktitle=pbteal!15,coltitle=black}
\else
\excludecomment{solution}
\fi
\newcounter{exo}
\usepackage{needspace}
\newenvironment{exo}[1]{\Needspace{8\baselineskip}\refstepcounter{exo}\medskip\noindent{\large\bfseries\color{pbblue}Exercise \theexo\ -- #1}\par\smallskip}{\par}
%% type of an exercise, printed after its title
\newcommand{\extype}[1]{{\normalfont\normalsize\itshape(#1)}}
%% heading of a block of exercises
\newcommand{\exoblock}[1]{\Needspace{10\baselineskip}\bigskip\noindent{\Large\bfseries\color{pbblue}#1}\par\nobreak\smallskip\noindent{\color{pbblue}\rule{\linewidth}{0.6pt}}\par\nobreak}

\begin{document}
\sloppy
\sheethead{TD -- PAC-Bayesian Learning\ifsol\ -- Solutions\fi}{This tutorial accompanies Sections~2--7 of the lecture notes
\emph{PAC-Bayesian Learning: a Theoretical Framework for Ensemble Methods}. Exercises~\ref{exo:quiz}--\ref{exo:testset}
are warm-ups (comprehension, classical short proofs, a first certificate by hand);
Exercises~\ref{exo:kl}--\ref{exo:selfbounding} follow the notes; Exercises~\ref{exo:votes}--\ref{exo:softmax} mix
applications and short proofs. Each title states the type of the exercise; those marked $(\star)$ are at Master level. Notation: $S\sim\D^m$,
voters $h:\X\to\{-1,+1\}$, posterior $Q$, prior $P$, Gibbs classifier $\GQ$, majority vote $\BQ$,
$W_Q(x,y)=\E_{h\sim Q}\ind{h(x)\neq y}$, joint error $e_Q=\E W_Q^2$, disagreement $d_Q$.
Unless stated otherwise, $\delta=0.05$, and $\klup(q,\varepsilon)=\sup\{p\in[q,1]:\kl(q\|p)\leq\varepsilon\}$.}

%% =====================================================================
%% =====================================================================
\exoblock{Part I -- Warm-up}

%% =====================================================================
\begin{exo}{True or false? \extype{comprehension}}\label{exo:quiz}
For each statement, say whether it is true or false and justify your answer in one or two sentences
(a short argument or a counterexample). $Q$ and $P$ are distributions on a set $\Hcal$ of voters with values in
$\{-1,+1\}$, and ties of the majority vote are counted as errors.
\begin{enumerate}[label=(\alph*)]
\item For every posterior $Q$, the majority vote is at least as good as the Gibbs classifier: $\RD(\BQ)\leq\RD(\GQ)$.
\item The empirical Gibbs risk $\RS(\GQ)$ always lies between the smallest and the largest empirical risk
      $\RS(h)$ of the voters $h$ in the support of $Q$.
\item The prior $P$ may depend on domain knowledge or on data \emph{other} than the sample $S$ on which the bound is
      computed, but not on $S$ itself.
\item For the bound to be valid, the posterior $Q$ must not depend on $S$ either.
\item Since a PAC-Bayesian bound holds for all posteriors simultaneously, we may compute it for a thousand posteriors
      and report the smallest value, without any correction.
\item ``With probability at least $0.95$'' means that at least $95\%$ of the voters $h\sim P$ satisfy $\RD(h)\leq$ the bound.
\item For the same empirical risk $\hat q$ and the same complexity term $\varepsilon$, the kl-based bound
      $\klup(\hat q,\varepsilon)$ is never larger than the Hoeffding-type bound $\hat q+\sqrt{\varepsilon/2}$, and the gain is
      largest when $\hat q$ is close to $0$.
\item The disagreement $d_Q$ can be estimated on unlabeled data, whereas the joint error $e_Q$ needs labels.
\item Taking $Q=P$ always gives the best bound, since it cancels the KL term.
\item The tandem bound $4e_Q$ is always tighter than the first-order bound $2\RD(\GQ)$.
\end{enumerate}
\end{exo}

\begin{solution}
\begin{enumerate}[label=(\alph*)]
\item \textbf{False.} Take $Q$ uniform on three voters and a distribution on which, on a fraction $r$ of the points,
      exactly two of the three voters err, and none elsewhere. Then $W_Q=\frac23$ on these points, the vote errs there,
      $\RD(\BQ)=r$, while $\RD(\GQ)=\E W_Q=\frac{2r}3<r$. The only general relation is $\RD(\BQ)\leq2\RD(\GQ)$
      (Exercise~\ref{exo:mv}); the vote is better only when the errors are spread over different points.
\item \textbf{True.} $\RS(\GQ)=\E_{h\sim Q}\RS(h)$ is an average of the $\RS(h)$ with weights $Q(h)$; an average lies
      between the minimum and the maximum of the averaged values. (This is \emph{not} true for the vote.)
\item \textbf{True.} In the proof, $P$ must be independent of $S$ to exchange $\E_{S}$ and $\E_{h\sim P}$ (Fubini); any
      information independent of $S$ is allowed, which is exactly what data-dependent priors learned on a separate
      sample exploit (Exercise~\ref{exo:ddprior}).
\item \textbf{False.} The bound holds, with probability $\geq1-\delta$, \emph{for all $Q$ simultaneously}; in
      particular for a posterior computed from $S$ in any way, even the minimizer of the bound (self-bounding learning,
      Exercise~\ref{exo:selfbounding}).
\item \textbf{True.} It is the same event (of probability $\geq1-\delta$ over $S$) that guarantees the bound for all
      posteriors; the minimum over a thousand of them is a data-dependent choice of $Q$, which is covered.
      Compare with (c): changing the \emph{prior} or $\delta$ changes the event and requires a union bound.
\item \textbf{False.} The probability is over the draw of the sample $S\sim\D^m$, not over voters or posteriors:
      for at least $95\%$ of the samples, the bound holds for \emph{every} posterior. Nothing is claimed about the
      proportion of good voters under $P$.
\item \textbf{True.} By Pinsker's inequality $\kl(\hat q\|p)\geq2(p-\hat q)^2$, every $p$ with $\kl(\hat q\|p)\leq\varepsilon$
      satisfies $p\leq\hat q+\sqrt{\varepsilon/2}$, so $\klup(\hat q,\varepsilon)\leq\hat q+\sqrt{\varepsilon/2}$. For $\hat q\to0$ the refined
      Pinsker inequality gives a width of order $\varepsilon$ instead of $\sqrt\varepsilon$ (Exercise~\ref{exo:testset}); for
      $\hat q$ close to $\frac12$ the two bounds nearly coincide.
\item \textbf{True.} $d_Q=\E_{h,h'\sim Q}\P_x(h(x)\neq h'(x))$ only compares predictions, whereas ``both $h$ and $h'$ err''
      requires knowing $y$.
\item \textbf{False.} $Q=P$ cancels the KL, but its empirical risk is the prior average $\E_P\RS(h)$, usually large.
      The bound trades the two terms, and its minimizer is typically a Gibbs posterior that moves away from $P$
      (Exercises~\ref{exo:dv} and~\ref{exo:gibbs4}).
\item \textbf{False.} Since $\RD(\GQ)=e_Q+\frac12d_Q$, we have $4e_Q\leq2\RD(\GQ)$ iff $d_Q\geq\RD(\GQ)$. For identical
      voters ($d_Q=0$), the tandem bound is $4\RD(\GQ)$, twice the first-order bound.
\end{enumerate}
\end{solution}

%% =====================================================================
\begin{exo}{From Markov to Hoeffding \extype{short proofs}}\label{exo:concentration}
In this exercise, $Z$ is a real random variable with mean $\mu$ and variance $v$, and $a>0$.
\begin{enumerate}
\item \emph{(Markov).} If $Z\geq0$, show that $\P(Z\geq a)\leq\frac{\E Z}a$ (hint: compare $a\ind{Z\geq a}$ and $Z$).
      Deduce that for every $\delta\in(0,1]$, $Z\leq\frac{\E Z}\delta$ with probability at least $1-\delta$. Give a
      distribution for which Markov's inequality is an equality.
\item \emph{(Chebyshev).} Show that $\P(|Z-\mu|\geq a)\leq\frac v{a^2}$.
\item \emph{(Cantelli).} For $t\geq0$, show that $\P(Z-\mu\geq a)\leq\frac{v+t^2}{(a+t)^2}$, then optimize over $t$ to obtain
      $\P(Z-\mu\geq a)\leq\frac{v}{v+a^2}$. Compare the two inequalities numerically for $v=0.01$ and $a=0.2$, and
      explain why the one-sided bound is always better than Chebyshev's.
\item \emph{(Union bound).} Show that $\P(\bigcup_{k=1}^KA_k)\leq\sum_{k=1}^K\P(A_k)$ for any events $A_1,\dots,A_K$.
      Deduce that if $K$ statements hold each with probability at least $1-\delta/K$, then they hold
      \emph{simultaneously} with probability at least $1-\delta$.
\item \emph{(Hoeffding's inequality).} Admit Hoeffding's lemma: if $\E X=0$ and $\alpha\leq X\leq\beta$, then
      $\E e^{sX}\leq e^{s^2(\beta-\alpha)^2/8}$ for all $s\in\R$. Let $Z_1,\dots,Z_m$ be i.i.d.\ in $[0,1]$ with mean
      $p$ and $\hat p=\frac1m\sum_iZ_i$.
      \begin{enumerate}[label=(\alph*)]
      \item Show that $\E e^{\lambda(p-\hat p)}\leq e^{\lambda^2/(8m)}$ for every $\lambda>0$.
      \item \emph{(Chernoff's method).} Apply Markov's inequality to $e^{\lambda(p-\hat p)}$ and optimize over $\lambda$ to
            show that $\P(p-\hat p\geq\varepsilon)\leq e^{-2m\varepsilon^2}$.
      \item Deduce that, with probability at least $1-\delta$, $p\leq\hat p+\sqrt{\ln(1/\delta)/(2m)}$, and give a two-sided
            version. Using question~4, give a bound valid simultaneously for $K$ fixed classifiers.
      \end{enumerate}
\item $(\star)$ \emph{(Hoeffding's lemma).} Let $\alpha\leq X\leq\beta$ with $\E X=0$ (so $\alpha<0<\beta$ unless $X=0$).
      Using the convexity of $x\mapsto e^{sx}$, show that $\E e^{sX}\leq e^{L(u)}$ with $u=s(\beta-\alpha)$,
      $\theta=\frac{-\alpha}{\beta-\alpha}$ and $L(u)=-\theta u+\ln(1-\theta+\theta e^u)$. Show that $L(0)=L'(0)=0$ and
      $L''\leq\frac14$, and conclude.
\end{enumerate}
\end{exo}

\begin{solution}
\begin{enumerate}
\item If $Z\geq a$, then $a\ind{Z\geq a}=a\leq Z$; otherwise $a\ind{Z\geq a}=0\leq Z$. Hence $a\ind{Z\geq a}\leq Z$ pointwise, and
      taking expectations, $a\P(Z\geq a)\leq\E Z$. With $a=\E Z/\delta$ (if $\E Z>0$; the case $\E Z=0$ is trivial since then
      $Z=0$ a.s.), $\P(Z\geq\E Z/\delta)\leq\delta$, i.e.\ $Z<\E Z/\delta$ with probability $\geq1-\delta$. \emph{Equality:} the
      pointwise inequality is an equality iff $Z\in\{0,a\}$; e.g.\ $Z=a$ with probability $r$ and $0$ otherwise gives
      $\P(Z\geq a)=r=\E Z/a$. This is exactly the worst case of the first-order bound ($W_Q\in\{0,\frac12\}$).
\item $\{|Z-\mu|\geq a\}=\{(Z-\mu)^2\geq a^2\}$, and Markov's inequality applied to the non-negative variable
      $(Z-\mu)^2$ gives $\P(|Z-\mu|\geq a)\leq\frac{\E(Z-\mu)^2}{a^2}=\frac v{a^2}$.
\item For $t\geq0$, if $Z-\mu\geq a$ then $Z-\mu+t\geq a+t>0$, hence $(Z-\mu+t)^2\geq(a+t)^2$. Markov applied to
      $(Z-\mu+t)^2$ gives
      \[\P(Z-\mu\geq a)\leq\frac{\E(Z-\mu+t)^2}{(a+t)^2}=\frac{v+2t\,\E(Z-\mu)+t^2}{(a+t)^2}=\frac{v+t^2}{(a+t)^2}=:g(t).\]
      Then $g'(t)=\frac{2t(a+t)-2(v+t^2)}{(a+t)^3}=\frac{2(at-v)}{(a+t)^3}$ vanishes at $t^\star=\frac va$, is negative before and
      positive after. Since $v+t^{\star2}=\frac{v(a^2+v)}{a^2}$ and $(a+t^\star)^2=\frac{(a^2+v)^2}{a^2}$, we get $g(t^\star)=\frac{v}{v+a^2}$.
      \emph{Numerically:} Chebyshev gives $\frac{0.01}{0.04}=0.25$ (for the two-sided event), Cantelli $\frac{0.01}{0.05}=0.20$
      (for the one-sided event). Since $v+a^2>a^2$, Cantelli is always smaller, and it is always $<1$; it is also
      the value $g(0)$ of Chebyshev's argument that is improved by shifting the variable: the free parameter $t$ is
      optimized. This is the mechanism behind the Chebyshev--Cantelli family and the C-bound (Exercise~\ref{exo:mv},
      question~4).
\item Pointwise, $\ind{\bigcup_kA_k}\leq\sum_k\ind{A_k}$: if the left-hand side equals $1$, at least one term of the
      sum equals $1$. Taking expectations gives the union bound. If $A_k$ is the event ``statement $k$ fails'', then
      $\P(A_k)\leq\delta/K$ and $\P(\text{some statement fails})\leq K\cdot\delta/K=\delta$.
\item (a) Write $p-\hat p=\sum_i\frac{p-Z_i}m$. The terms are independent, centred, with values in $[\frac{p-1}m,\frac pm]$, an
      interval of length $\frac1m$. By independence and Hoeffding's lemma,
      \[\E e^{\lambda(p-\hat p)}=\prod_{i=1}^m\E e^{\lambda\frac{p-Z_i}m}\leq\prod_{i=1}^me^{\lambda^2/(8m^2)}=e^{\lambda^2/(8m)}.\]
      (b) For $\lambda>0$, $x\mapsto e^{\lambda x}$ is increasing, so by Markov
      $\P(p-\hat p\geq\varepsilon)=\P(e^{\lambda(p-\hat p)}\geq e^{\lambda\varepsilon})\leq e^{-\lambda\varepsilon}\E e^{\lambda(p-\hat p)}\leq e^{-\lambda\varepsilon+\lambda^2/(8m)}$.
      The exponent is minimal at $\lambda=4m\varepsilon$, where it equals $-4m\varepsilon^2+2m\varepsilon^2=-2m\varepsilon^2$.
      (c) Setting $e^{-2m\varepsilon^2}=\delta$ gives $\varepsilon=\sqrt{\ln(1/\delta)/(2m)}$: with probability $\geq1-\delta$,
      $p<\hat p+\varepsilon$. The same argument applied to $1-Z_i$ bounds $\P(\hat p-p\geq\varepsilon)$; with the union bound
      (each side at $\delta/2$), $|p-\hat p|\leq\sqrt{\ln(2/\delta)/(2m)}$ with probability $\geq1-\delta$. For $K$ classifiers
      $h_1,\dots,h_K$ fixed \emph{before} seeing $S$, applying the one-sided bound to each with $\delta/K$ gives,
      with probability $\geq1-\delta$, for all $k$: $\RD(h_k)\leq\RS(h_k)+\sqrt{\ln(K/\delta)/(2m)}$. The price of $K$ classifiers
      is only logarithmic.
\item By convexity, for $x\in[\alpha,\beta]$, $e^{sx}\leq\frac{\beta-x}{\beta-\alpha}e^{s\alpha}+\frac{x-\alpha}{\beta-\alpha}e^{s\beta}$ (the
      chord lies above the graph). Taking expectations with $\E X=0$:
      $\E e^{sX}\leq\frac{\beta e^{s\alpha}-\alpha e^{s\beta}}{\beta-\alpha}=(1-\theta)e^{s\alpha}+\theta e^{s\beta}$, with $\theta=\frac{-\alpha}{\beta-\alpha}\in[0,1]$.
      Since $s\alpha=-\theta u$ and $s\beta=(1-\theta)u$, this equals $e^{-\theta u}(1-\theta+\theta e^u)=e^{L(u)}$. Now $L(0)=\ln1=0$,
      $L'(u)=-\theta+\frac{\theta e^u}{1-\theta+\theta e^u}$ so $L'(0)=0$, and $L''(u)=\frac{\theta e^u(1-\theta)}{(1-\theta+\theta e^u)^2}=r(1-r)\leq\frac14$
      with $r=\frac{\theta e^u}{1-\theta+\theta e^u}\in[0,1]$. By Taylor's formula with Lagrange remainder,
      $L(u)=L(0)+uL'(0)+\frac{u^2}2L''(\xi)\leq\frac{u^2}8=\frac{s^2(\beta-\alpha)^2}8$.
\end{enumerate}
\emph{What we learn:} every concentration inequality of the course is Markov's inequality applied to a well-chosen
non-negative function ($Z$, $(Z-\mu)^2$, $(Z-\mu+t)^2$, $e^{\lambda(p-\hat p)}$), followed by the optimization of a free
parameter; PAC-Bayesian bounds follow the same pattern, with the change of measure in the middle.
\end{solution}

%% =====================================================================
\begin{exo}{Certifying one classifier \extype{application}}\label{exo:testset}
A classifier $f$ (for instance a random forest) has been learned on a sample $S_1$. It is then evaluated on a test
sample $T$ of $n=200$ examples drawn i.i.d.\ from $\D$, independently of $S_1$, and makes $10$ errors.
\begin{enumerate}
\item Why is $n\widehat R_T(f)$ a binomial variable $\mathrm{Bin}(n,\RD(f))$ conditionally on $S_1$? Which assumption
      would break if $T$ had been used to tune the hyperparameters of $f$?
\item Compute the Hoeffding test-set bound $\widehat R_T(f)+\sqrt{\ln(1/\delta)/(2n)}$.
\item The kl test-set bound is $\klup\big(\widehat R_T(f),\frac{\ln(1/\delta)}n\big)$. Compute $\varepsilon=\ln(1/\delta)/n$,
      check that $\kl(0.05\|0.15)>\varepsilon$, and perform four steps of bisection on $[0.05,0.15]$ (compute
      $\kl(0.05\|c)$ at each midpoint $c$). Compare the result with the Pinsker relaxation $\hat q+\sqrt{\varepsilon/2}$ and
      with the refined relaxation $\hat q+\sqrt{2\hat q\varepsilon}+2\varepsilon$.
\item In fact, the engineer trained $K=20$ classifiers (different hyperparameters), evaluated all of them on $T$ and
      kept the one with the fewest test errors, which is $f$. Why are the certificates of questions~2 and~3 no longer
      valid? Correct them with a union bound, for $K=20$ and $K=1000$ (for the kl bound you may use
      \texttt{kl\_inv\_upper} of \texttt{pbtools}).
\item What test size $n$ guarantees a Hoeffding width $\sqrt{\ln(K/\delta)/(2n)}\leq0.02$, for $K\in\{1,20,1000\}$? Comment.
\item Show that the corrected bound of question~4 is the Occam bound (Section~2.4 of the notes) for a well-chosen
      prior, and give its KL term.
\end{enumerate}
\end{exo}

\begin{solution}
\begin{enumerate}
\item Conditionally on $S_1$, $f$ is a fixed function. The test examples are i.i.d.\ and independent of $S_1$, so the
      indicators $\ind{f(x_i)\neq y_i}$, $i\leq n$, are i.i.d.\ Bernoulli variables with parameter
      $\P_\D(f(x)\neq y)=\RD(f)$, and their sum is binomial. If $T$ had been used to tune $f$, then $f$ would depend on
      $T$: the indicators would no longer be independent Bernoulli$(\RD(f))$ variables (the tuning selects a
      classifier that happens to do well on $T$), and $\widehat R_T(f)$ would be optimistically biased.
\item $\widehat R_T(f)=\frac{10}{200}=0.05$ and $\ln(1/\delta)=\ln20=2.996$, so the bound is
      $0.05+\sqrt{2.996/400}=0.05+0.0865=0.1365$.
\item $\varepsilon=\frac{2.996}{200}=0.01498$. At the right end, $\kl(0.05\|0.15)=0.0507>\varepsilon$, so
      $\klup(0.05,\varepsilon)\in[0.05,0.15]$ (the function $p\mapsto\kl(0.05\|p)$ is increasing on $[0.05,1]$). Bisection:
      \begin{center}
      \begin{tabular}{cccc}\toprule
      step & midpoint $c$ & $\kl(0.05\|c)$ & new interval\\\midrule
      1 & $0.1$ & $0.0167>\varepsilon$ & $[0.05,0.1]$\\
      2 & $0.075$ & $0.0051\leq\varepsilon$ & $[0.075,0.1]$\\
      3 & $0.0875$ & $0.0103\leq\varepsilon$ & $[0.0875,0.1]$\\
      4 & $0.09375$ & $0.0134\leq\varepsilon$ & $[0.09375,0.1]$\\\bottomrule
      \end{tabular}
      \end{center}
      After four steps we know that $\klup\in[0.094,0.1]$, and we may already report the certified value $0.1$ (the
      right end is always a valid upper bound). One more step ($\kl(0.05\|0.0969)=0.01500>\varepsilon$) and further
      iterations give $\klup(0.05,0.01498)=0.0968$. The Pinsker relaxation is $0.05+\sqrt{0.00749}=0.1365$, i.e.\ the
      Hoeffding bound of question~2, and the refined relaxation is $0.05+\sqrt{2\times0.05\times0.01498}+2\times0.01498=0.05+0.0387+0.0300=0.1187$.
      Hence $0.0968<0.1187<0.1365$: for a small empirical risk, the kl bound is much tighter (width $0.047$ instead of $0.087$).
\item Each of the $K$ certificates is valid for a classifier fixed before looking at $T$; but the classifier kept
      is chosen \emph{with} $T$, among $K$ candidates, precisely because its test error is small. The event on which
      its certificate holds may be the one event that fails (a winner's curse). With the union bound (question~4 of
      Exercise~\ref{exo:concentration}), applying each certificate with $\delta/K$ makes all $K$ of them hold
      simultaneously with probability $\geq1-\delta$, and in particular the one of the selected classifier. For $K=20$:
      $\ln(K/\delta)=\ln400=5.99$, Hoeffding gives $0.05+\sqrt{5.99/400}=0.1724$ and kl gives
      $\klup(0.05,0.0300)=0.1215$. For $K=1000$: $\ln(20000)=9.90$, Hoeffding $0.05+\sqrt{9.90/400}=0.2073$, kl
      $\klup(0.05,0.0495)=0.1484$.
\item $n\geq\frac{\ln(K/\delta)}{2\times0.02^2}=1250\ln(K/\delta)$, i.e.\ $n\geq3745$ for $K=1$, $n\geq7490$ for $K=20$ and
      $n\geq12380$ for $K=1000$. The precision improves like $1/\sqrt n$ (four times more examples to halve the
      width), while the number of candidates only enters through $\ln K$: comparing $1000$ models ``only'' costs
      about $3.3$ times more test data than certifying a single one.
\item Take the uniform prior $P=\frac1K\sum_k\delta_{f_k}$ on the $K$ candidates and the Dirac posterior on the selected one.
      Then $\KL(\delta_f\|P)=\ln\frac1{P(f)}=\ln K$, and the Occam bound $\RD(f)\leq\klup\big(\widehat R_T(f),\frac{\ln\frac1{P(f)}+\ln\frac1\delta}n\big)$
      is exactly the corrected kl bound: the union bound is the simplest PAC-Bayesian bound, and $\ln K$ is its
      complexity term (Exercise~\ref{exo:occam}).
\end{enumerate}
\emph{What we learn:} a test-set bound is only valid for a predictor chosen independently of the test sample;
every data-dependent selection has a price, here $\ln K$ nats, and the kl inversion is the right tool for small risks.
\end{solution}

%% =====================================================================
%% =====================================================================
\exoblock{Part II -- Core exercises}

%% =====================================================================
\begin{exo}{The Kullback--Leibler divergence \extype{short proofs}}\label{exo:kl}
Let $\Hcal=\{h_1,\dots,h_n\}$ be finite and $P,Q$ two distributions on $\Hcal$ with $P_i>0$.
\begin{enumerate}
\item Using Jensen's inequality, show that $\KL(Q\|P)\geq0$, with equality iff $Q=P$.
\item Assume $P$ uniform. Show that $\KL(Q\|P)=\ln n-H(Q)$ where $H(Q)=-\sum_iQ_i\ln Q_i$. What are the
      minimal and maximal values of $\KL(Q\|P)$? Compute it when $Q$ is uniform over $k\leq n$ voters.
\item Show that $\KL(Q\otimes Q\,\|\,P\otimes P)=2\KL(Q\|P)$.
\item Show that $\KL\big(\mathcal N(\mu_1,\sigma^2I_d)\,\|\,\mathcal N(\mu_0,\sigma^2I_d)\big)=\frac{\|\mu_1-\mu_0\|^2}{2\sigma^2}$.
\item $(\star)$ Binary kl. Fix $q\in[0,1]$. Show that $p\mapsto\kl(q\|p)$ is convex on $(0,1)$, minimal at $p=q$, and
      prove Pinsker's inequality $\kl(q\|p)\geq2(q-p)^2$ by studying $g(p)=\kl(q\|p)-2(q-p)^2$.
\end{enumerate}
\end{exo}

\begin{solution}
\begin{enumerate}
\item \emph{Idea: $-\KL$ is the expectation of a logarithm, and $\ln$ is concave.} Let $I=\{i:Q_i>0\}$ be the
      support of $Q$. By definition, $-\KL(Q\|P)=\sum_{i\in I}Q_i\ln\frac{P_i}{Q_i}$. Since the $Q_i$, $i\in I$, are
      weights summing to one and $\ln$ is concave, Jensen's inequality gives
      \[-\KL(Q\|P)\leq\ln\Big(\sum_{i\in I}Q_i\frac{P_i}{Q_i}\Big)=\ln\Big(\sum_{i\in I}P_i\Big)\leq\ln1=0.\]
      \emph{Equality case.} Jensen's inequality for a strictly concave function is an equality iff the variable is
      constant: $P_i/Q_i=c$ for all $i\in I$. The second inequality is an equality iff $\sum_{i\in I}P_i=1$, i.e.\
      $I=\Hcal$ (since all $P_i>0$). Then $c=\sum_iP_i/\sum_iQ_i=1$ and $Q=P$. Conversely $\KL(P\|P)=0$.
\item With $P_i=\frac1n$: $\KL(Q\|P)=\sum_iQ_i\ln(nQ_i)=\ln n\sum_iQ_i+\sum_iQ_i\ln Q_i=\ln n-H(Q)$.
      The entropy satisfies $H(Q)\geq0$ (each term $-Q_i\ln Q_i\geq0$), with equality for a Dirac, and
      $H(Q)\leq\ln n$ (this is question~1 applied to $\KL(Q\|\text{uniform})\geq0$), with equality for the uniform
      distribution. Hence $0\leq\KL(Q\|P)\leq\ln n$: the KL is $0$ when the posterior is the prior and $\ln n$ when it
      selects a single voter (the Occam price). For $Q$ uniform on $k$ voters, $H(Q)=k\cdot\frac1k\ln k=\ln k$ and
      $\KL=\ln\frac nk$: selecting a tenth of the voters costs $\ln10\approx2.3$ nats, whatever $n$.
\item Write $(Q\otimes Q)_{ij}=Q_iQ_j$. Then $\ln\frac{Q_iQ_j}{P_iP_j}=\ln\frac{Q_i}{P_i}+\ln\frac{Q_j}{P_j}$, so
      \begin{align*}\KL(Q^2\|P^2)&=\sum_{i,j}Q_iQ_j\ln\frac{Q_i}{P_i}+\sum_{i,j}Q_iQ_j\ln\frac{Q_j}{P_j}\\
      &=\sum_iQ_i\ln\frac{Q_i}{P_i}+\sum_jQ_j\ln\frac{Q_j}{P_j}=2\KL(Q\|P),\end{align*}
      using $\sum_jQ_j=1$ in the first sum and $\sum_iQ_i=1$ in the second. This is why second-order bounds
      (tandem, disagreement) pay a doubled KL.
\item The two densities have the same normalizing constant, so the log-ratio is
      $\ln\frac{q(\theta)}{p(\theta)}=\frac{\|\theta-\mu_0\|^2-\|\theta-\mu_1\|^2}{2\sigma^2}$. Write
      $\theta=\mu_1+\sigma\varepsilon$ with $\varepsilon\sim\mathcal N(0,I_d)$ under $Q$. Then
      $\|\theta-\mu_1\|^2=\sigma^2\|\varepsilon\|^2$ and
      $\|\theta-\mu_0\|^2=\|\mu_1-\mu_0\|^2+2\sigma\langle\varepsilon,\mu_1-\mu_0\rangle+\sigma^2\|\varepsilon\|^2$. The
      terms in $\|\varepsilon\|^2$ cancel, and $\E\langle\varepsilon,\mu_1-\mu_0\rangle=0$, so
      $\KL=\E_Q\ln\frac qp=\frac{\|\mu_1-\mu_0\|^2}{2\sigma^2}$. Moving the centre of a Gaussian posterior costs the
      squared distance travelled, measured in units of the variance.
\item For $p\in(0,1)$, $\partial_p\kl(q\|p)=-\frac qp+\frac{1-q}{1-p}=\frac{-q(1-p)+p(1-q)}{p(1-p)}=\frac{p-q}{p(1-p)}$,
      which is negative for $p<q$ and positive for $p>q$: the function decreases then increases, and is minimal at
      $p=q$ where it vanishes. Moreover $\partial_p^2\kl(q\|p)=\frac q{p^2}+\frac{1-q}{(1-p)^2}>0$ (unless $q\in\{0,1\}$,
      where it is still $\geq0$): the function is convex.
      \emph{Pinsker.} Let $g(p)=\kl(q\|p)-2(q-p)^2$. Then $g(q)=0$ and
      \[g'(p)=\frac{p-q}{p(1-p)}-4(p-q)=(p-q)\Big(\frac{1}{p(1-p)}-4\Big).\]
      Since $p(1-p)\leq\frac14$, the bracket is non-negative, so $g'$ has the sign of $p-q$: $g$ decreases on $(0,q)$,
      increases on $(q,1)$, and its minimum is $g(q)=0$. Hence $\kl(q\|p)\geq2(q-p)^2$. The inequality is tight
      when $p(1-p)\approx\frac14$, i.e.\ around $p=\frac12$, which is why the kl and Hoeffding bounds coincide there.
\end{enumerate}
\end{solution}

%% =====================================================================
\begin{exo}{Change of measure and Gibbs posterior \extype{short proofs}}\label{exo:dv}
Let $\phi:\Hcal\to\R$ and $P$ a distribution on a finite set $\Hcal$. Define $Q_\phi(h)=\frac{P(h)e^{\phi(h)}}{\E_{P}e^{\phi}}$.
\begin{enumerate}
\item Show that for any $Q$, $\KL(Q\|Q_\phi)=\KL(Q\|P)-\E_Q\phi+\ln\E_Pe^\phi$.
\item Deduce the change of measure inequality $\E_Q\phi\leq\KL(Q\|P)+\ln\E_Pe^\phi$ and the variational formula
      $\ln\E_Pe^{\phi}=\max_Q\{\E_Q\phi-\KL(Q\|P)\}$.
\item Let $\lambda>0$. Show that the minimizer of $Q\mapsto\RS(\GQ)+\frac{1}{\lambda m}\KL(Q\|P)$ is the Gibbs posterior
      $Q_\lambda(h)\propto P(h)e^{-\lambda m\RS(h)}$, and give the minimal value.
\item Describe $Q_\lambda$ when $\lambda\to0$ and $\lambda\to+\infty$. Interpret $\lambda$.
\item Recover $Q_\lambda$ directly with a Lagrange multiplier for the constraint $\sum_hQ(h)=1$.
\end{enumerate}
\end{exo}

\begin{solution}
\begin{enumerate}
\item By definition of $Q_\phi$, $\ln\frac{Q(h)}{Q_\phi(h)}=\ln\frac{Q(h)}{P(h)}-\phi(h)+\ln\E_Pe^\phi$ (the last term does
      not depend on $h$). Taking the expectation under $Q$:
      $\KL(Q\|Q_\phi)=\KL(Q\|P)-\E_Q\phi+\ln\E_Pe^\phi$.
\item By Exercise~\ref{exo:kl}, $\KL(Q\|Q_\phi)\geq0$, hence $\E_Q\phi\leq\KL(Q\|P)+\ln\E_Pe^\phi$ for every $Q$: this is
      the change of measure inequality. It is an equality iff $\KL(Q\|Q_\phi)=0$, i.e.\ iff $Q=Q_\phi$. Therefore
      $\E_Q\phi-\KL(Q\|P)\leq\ln\E_Pe^\phi$ for all $Q$, with equality at $Q_\phi$: the maximum over $Q$ is reached
      and equals $\ln\E_Pe^\phi$. \emph{Interpretation:} the left-hand side of the change of measure concerns a
      posterior that may depend on the data, the right-hand side only involves the prior; the price of the
      transfer is exactly the KL.
\item Apply question~2 to $\phi=-\lambda m\RS$ (a function of $h$ once $S$ is fixed):
      $-\lambda m\,\E_Q\RS-\KL(Q\|P)\leq\ln\E_Pe^{-\lambda m\RS}$, i.e., dividing by $-\lambda m<0$,
      \[\RS(\GQ)+\frac{\KL(Q\|P)}{\lambda m}\geq-\frac1{\lambda m}\ln\E_{h\sim P}e^{-\lambda m\RS(h)},\]
      with equality iff $Q=Q_\phi$, i.e.\ $Q(h)\propto P(h)e^{-\lambda m\RS(h)}$. The minimal value is the right-hand side,
      a ``soft minimum'' of the empirical risks (it tends to $\min_h\RS(h)$ when $\lambda\to\infty$ on a finite class).
\item When $\lambda\to0$, $e^{-\lambda m\RS(h)}\to1$ for every $h$ and $Q_\lambda\to P$: the data are ignored.
      When $\lambda\to\infty$, divide numerator and denominator by $e^{-\lambda m r^\star}$ with $r^\star=\min_h\RS(h)$:
      the weight of every $h$ with $\RS(h)>r^\star$ tends to $0$, and $Q_\lambda$ converges to $P$ restricted to the
      empirical risk minimizers and renormalized. $\lambda$ is an inverse temperature: it sets how much a
      difference of one training error changes the relative weights (by a factor $e^{\lambda}$).
\item The objective $F(Q)=\sum_hQ(h)\RS(h)+\frac1{\lambda m}\sum_hQ(h)\ln\frac{Q(h)}{P(h)}$ is strictly convex
      ($x\mapsto x\ln x$ is strictly convex). With the Lagrangian $\mathcal L(Q,\nu)=F(Q)+\nu(\sum_hQ(h)-1)$,
      \[\frac{\partial\mathcal L}{\partial Q(h)}=\RS(h)+\frac{1}{\lambda m}\Big(\ln\frac{Q(h)}{P(h)}+1\Big)+\nu=0
      \iff Q(h)=P(h)\,e^{-\lambda m\RS(h)}\,e^{-1-\lambda m\nu}.\]
      The constant $e^{-1-\lambda m\nu}$ is fixed by $\sum_hQ(h)=1$, which gives the Gibbs posterior. The positivity
      constraints are automatically satisfied, and by strict convexity this stationary point is the unique minimum.
\end{enumerate}
\end{solution}

%% =====================================================================
\begin{exo}{From the general theorem to classical bounds \extype{proofs and application}}\label{exo:bounds}
Recall the general theorem: for convex $\Delta$, w.p.\ $\geq1-\delta$, for all $Q$,
$\Delta(\RS(\GQ),\RD(\GQ))\leq\frac1m[\KL(Q\|P)+\ln\frac{1}{\delta}\E_{S'}\E_{h\sim P}e^{m\Delta(\widehat R_{S'}(h),\RD(h))}]$.
\begin{enumerate}
\item Explain precisely where each of the following is used in the proof: convexity of $\Delta$, the change of measure,
      Markov's inequality, the independence of $P$ from $S$.
\item \emph{(Hoeffding-type bound).} Take $\Delta(q,p)=\lambda(p-q)$ for a fixed $\lambda>0$. Using Hoeffding's lemma, show that
      w.p.\ $\geq1-\delta$, for all $Q$: $\RD(\GQ)\leq\RS(\GQ)+\frac{\KL(Q\|P)+\ln\frac1\delta}{\lambda m}+\frac\lambda8$.
\item Minimize the right-hand side in $\lambda$. Why is it not allowed to plug this value into the bound? Propose a
      correct way to optimize $\lambda$ and give the resulting price.
\item \emph{(Catoni's bound).} Let $C>0$ and $\mathcal F_C(p)=-\ln(1-p(1-e^{-C}))$. Show that for a fixed voter with risk $p$ and
      the 0-1 loss, $\E_{S'}e^{-Cm\widehat R_{S'}(h)}=(1-p(1-e^{-C}))^m$. Deduce Catoni's bound.
\item Numerical application: $m=1000$, $\delta=0.05$, $\RS(\GQ)=0.1$, $\KL(Q\|P)=5$. Compute McAllester's bound
      $\RS+\sqrt{(\KL+\ln\frac{2\sqrt m}\delta)/(2m)}$ and the Hoeffding-type bound for $\lambda\in\{0.1,0.2\}$.
\end{enumerate}
\end{exo}

\begin{solution}
\begin{enumerate}
\item The proof has three steps (Section 3.2 of the notes).
      \emph{Convexity of $\Delta$} is used in the first step (Jensen): the risks of $\GQ$ are averages over $Q$ of the
      risks of the voters, and $\Delta(\E_Q\RS(h),\E_Q\RD(h))\leq\E_Q\Delta(\RS(h),\RD(h))$. This reduces the problem to
      single voters.
      \emph{The change of measure} is the second step: $\E_Q[m\Delta]\leq\KL(Q\|P)+\ln\E_Pe^{m\Delta}$. It removes $Q$
      from the random part of the inequality; since it holds for all $Q$ at once, so will the final bound.
      \emph{Markov's inequality} is the third and only probabilistic step: the non-negative variable
      $Z(S)=\E_Pe^{m\Delta(\RS(h),\RD(h))}$ satisfies $Z(S)\leq\E_{S'}Z(S')/\delta$ with probability $\geq1-\delta$.
      \emph{Independence of $P$ and $S$} is used inside this step, to exchange $\E_{S'}$ and $\E_{h\sim P}$ (Fubini):
      for each fixed $h$, $\widehat R_{S'}(h)$ is an average of i.i.d.\ variables, whose exponential moment can be
      computed. If $P$ depended on $S'$, this exchange would be meaningless.
\item Here $\Delta(q,p)=\lambda(p-q)$ is linear, hence convex. For a fixed $h$, $\widehat R_{S'}(h)=\frac1m\sum_i\ell_i$ with
      i.i.d.\ $\ell_i\in[0,1]$ of mean $\RD(h)$, so by independence
      \[\E_{S'}e^{m\lambda(\RD(h)-\widehat R_{S'}(h))}=\prod_{i=1}^m\E\,e^{\lambda(\RD(h)-\ell_i)}\leq\big(e^{\lambda^2/8}\big)^m,\]
      by Hoeffding's lemma (a centred variable with range $1$ has $\E e^{sX}\leq e^{s^2/8}$). Hence
      $\mathcal I_\Delta(m)\leq e^{m\lambda^2/8}$, and the theorem gives
      $\lambda(\RD(\GQ)-\RS(\GQ))\leq\frac1m\big[\KL(Q\|P)+\ln\frac1\delta\big]+\frac{\lambda^2}8$. Dividing by $\lambda$ gives the result.
\item Let $A=\frac{\KL(Q\|P)+\ln(1/\delta)}{m}$. The function $\lambda\mapsto\frac A\lambda+\frac\lambda8$ is minimal at
      $\lambda^\star=\sqrt{8A}$, where it equals $\sqrt{A/2}$: we recover a McAllester-type bound
      $\RS(\GQ)+\sqrt{\frac{\KL+\ln(1/\delta)}{2m}}$. \emph{Why it is forbidden:} $\lambda$ is a parameter of $\Delta$, i.e.\ of
      the probabilistic event on which the bound holds; this event must be fixed before drawing $S$. But
      $\lambda^\star$ depends on $\KL(Q\|P)$, hence on $Q$ and on $S$. Using it amounts to choosing, after seeing the
      data, among infinitely many events, each of probability $\geq1-\delta$: nothing guarantees that the chosen one
      holds. \emph{Correct way:} fix a grid $\lambda_1<\dots<\lambda_k$ in advance, apply the bound to each $\lambda_j$ with
      confidence $\delta/k$ (union bound): with probability $\geq1-\delta$ all $k$ bounds hold, and we can take the best
      one. The price is $\frac{\ln k}{\lambda m}$ (i.e.\ $\ln k$ added to $\ln\frac1\delta$) plus the error due to the
      discretization; a geometric grid with ratio $2$ covering $[10^{-3},1]$ has $k=10$ points and costs $\ln10\approx2.3$ nats.
\item For the 0-1 loss, $m\widehat R_{S'}(h)\sim\mathrm{Bin}(m,p)$ with $p=\RD(h)$, whose moment generating function is
      $\E e^{sB}=(1-p+pe^s)^m$. At $s=-C$: $\E_{S'}e^{-Cm\widehat R_{S'}(h)}=(1-p+pe^{-C})^m=(1-p(1-e^{-C}))^m=e^{-m\mathcal F_C(p)}$.
      Take $\Delta(q,p)=\mathcal F_C(p)-Cq$ (convex: $\mathcal F_C$ is convex, being $-\ln$ of an affine decreasing function,
      and $-Cq$ is linear). Then $\E_{S'}e^{m\Delta(\widehat R_{S'}(h),p)}=e^{m\mathcal F_C(p)}e^{-m\mathcal F_C(p)}=1$, so
      $\mathcal I_\Delta(m)=1$ and, with probability $\geq1-\delta$, for all $Q$:
      $\mathcal F_C(\RD(\GQ))\leq C\RS(\GQ)+\frac{\KL(Q\|P)+\ln(1/\delta)}m$. Since $\mathcal F_C$ is increasing with inverse
      $\mathcal F_C^{-1}(u)=\frac{1-e^{-u}}{1-e^{-C}}$, we obtain
      $\RD(\GQ)\leq\frac{1}{1-e^{-C}}\big[1-\exp\big(-C\RS(\GQ)-\frac{\KL+\ln(1/\delta)}m\big)\big]$. (For losses in $[0,1]$
      that are not binary, the equality becomes an inequality by convexity of $u\mapsto e^{-Cu}$.)
\item $\ln\frac{2\sqrt{1000}}{0.05}=\ln1264.9\approx7.14$, so McAllester's bound is
      $0.1+\sqrt{\frac{5+7.14}{2000}}=0.1+0.078=0.178$. For the Hoeffding-type bound, $\ln\frac1{0.05}\approx3.00$ and
      $\KL+\ln\frac1\delta=8.00$. With $\lambda=0.1$: $0.1+\frac{8.00}{100}+\frac{0.1}8=0.1+0.080+0.0125=0.1925$. With
      $\lambda=0.2$: $0.1+\frac{8.00}{200}+0.025=0.165$. The optimal value would be $\lambda^\star=\sqrt{8\times0.008}=0.253$,
      giving $0.1+\sqrt{0.004}=0.163$. Remember that $\lambda=0.2$ is legitimate only if it was chosen before seeing the
      data; if it was selected among $\{0.1,0.2\}$ after computing both bounds, each must be computed with $\delta/2$,
      which gives $0.1+\frac{5+\ln40}{200}+0.025=0.1+0.0434+0.025\approx0.168$ for $\lambda=0.2$.
\end{enumerate}
\end{solution}

%% =====================================================================
\begin{exo}{Occam's razor as a PAC-Bayesian bound \extype{comprehension}}\label{exo:occam}
\begin{enumerate}
\item Let $\Hcal$ be countable and $Q=\delta_{h}$ a Dirac on $h\in\Hcal$. Compute $\KL(\delta_h\|P)$. What does Seeger's bound give
      for a single voter? Compare with the Occam bound of the lecture notes.
\item Why can we not use Dirac posteriors when $\Hcal$ is uncountable (e.g.\ linear classifiers with a Gaussian prior)?
\item Let $\Hcal=\{h_1,h_2,\dots\}$ be countable and $P(h_k)=\frac{6}{\pi^2k^2}$. Check that $P$ is a probability and give
      the complexity term of $h_k$. Interpret this prior in terms of \emph{description length}.
\end{enumerate}
\end{exo}

\begin{solution}
\begin{enumerate}
\item $\KL(\delta_h\|P)=\sum_{h'}\delta_h(h')\ln\frac{\delta_h(h')}{P(h')}=1\cdot\ln\frac1{P(h)}$. Seeger's bound with
      $Q=\delta_h$ (for which $\RS(G_{\delta_h})=\RS(h)$) gives, simultaneously for all $h$,
      \[\RD(h)\leq\klup\Big(\RS(h),\frac{\ln\frac1{P(h)}+\ln\frac{2\sqrt m}\delta}{m}\Big).\]
      The Occam bound of the notes has $\ln\frac1\delta$ instead of $\ln\frac{2\sqrt m}\delta$: it is slightly tighter
      because it only covers Dirac posteriors (a union bound over countably many events), whereas the PAC-Bayesian
      bound covers all posteriors simultaneously. For $m=1000$ the difference is $\ln(2\sqrt{1000})\approx4.1$ nats.
\item If $P$ has a density (e.g.\ a Gaussian on $\R^d$), then $P(\{h\})=0$ for every $h$, and $\delta_h$ is not
      absolutely continuous with respect to $P$: $\KL(\delta_h\|P)=+\infty$ and the bound is void. One must use
      ``spread'' posteriors, e.g.\ $\mathcal N(\mathbf w,\sigma^2I)$ around the chosen classifier; the bound then
      certifies the Gibbs classifier around $\mathbf w$, and the margin or the majority vote relates it to $\mathbf w$
      itself (Exercise~\ref{exo:linear}).
\item $\sum_{k\geq1}\frac1{k^2}=\frac{\pi^2}6$, so $\sum_kP(h_k)=1$ and $P$ is a probability. The complexity of $h_k$ is
      $\ln\frac1{P(h_k)}=2\ln k+\ln\frac{\pi^2}6\approx2\ln k+0.50$. \emph{Description length:} by Kraft's inequality, any
      prefix code assigning $\ell(h)$ nats to $h$ defines a sub-probability $P(h)=e^{-\ell(h)}$, and conversely. The
      prior thus encodes a way of describing hypotheses; the bound says that a hypothesis is cheap if it has a short
      description \emph{chosen before seeing the data}. Here the index $k$ is described with about $2\ln k$ nats, so
      hypotheses listed early are favoured: this is Occam's razor made quantitative.
\end{enumerate}
\end{solution}

%% =====================================================================
\begin{exo}{Bounds on the risk of the majority vote \extype{proofs and application}}\label{exo:mv}
We assume voters in $\{-1,+1\}$, 0-1 loss, and count ties as errors.
\begin{enumerate}
\item Show that $\RD(\BQ)=\P(W_Q\geq\frac12)$ and $\RD(\GQ)=\E W_Q$. Deduce $\RD(\BQ)\leq2\RD(\GQ)$. Give an
      example of distribution of $W_Q$ for which this inequality is an equality.
\item Show that $e_Q=\E[W_Q^2]$, $d_Q=\E[2W_Q(1-W_Q)]$, and $\RD(\GQ)=e_Q+\frac12d_Q$. Show also that
      $d_Q\leq2\RD(\GQ)(1-\RD(\GQ))$.
\item Prove the tandem bound $\RD(\BQ)\leq4e_Q$. Show that it is better than the first-order bound iff $d_Q\geq\RD(\GQ)$.
\item For $\mu<\frac12$, show that $\RD(\BQ)\leq\frac{\E(W_Q-\mu)^2}{(1/2-\mu)^2}$. Find the optimal $\mu^\star$ and show that
      the optimal value is $\frac{\Var W_Q}{\Var W_Q+(1/2-\E W_Q)^2}$ (Cantelli's inequality).
\item Deduce the C-bound $\RD(\BQ)\leq1-\frac{(1-2\RD(\GQ))^2}{1-2d_Q}$ (when $\RD(\GQ)<\frac12$).
\item \emph{(Condorcet).} $Q$ is uniform over $n$ voters erring independently with probability $p$. Compute $\RD(\GQ)$,
      $e_Q$, $d_Q$ as functions of $n$ and $p$. For $p=0.3$, compute the exact risk of the vote and the three bounds
      for $n=3$, then the C-bound for $n=25$ and its limit when $n\to\infty$. Comment.
\end{enumerate}
\end{exo}

\begin{solution}
\begin{enumerate}
\item For a fixed $(x,y)$, $y\,\E_Qh(x)=\E_Q[yh(x)]=\E_Q[1-2\ind{h(x)\neq y}]=1-2W_Q(x,y)$, because $yh(x)=1$ when
      $h$ is right and $-1$ when it errs. The vote errs (ties counted as errors) iff $y\,\E_Qh(x)\leq0$, i.e.\ iff
      $W_Q(x,y)\geq\frac12$. Taking the probability over $(x,y)\sim\D$: $\RD(\BQ)=\P(W_Q\geq\frac12)$. By Fubini,
      $\E_\D W_Q=\E_Q\E_\D\ind{h(x)\neq y}=\E_Q\RD(h)=\RD(\GQ)$. Markov's inequality applied to $W_Q\geq0$ gives
      $\P(W_Q\geq\frac12)\leq\frac{\E W_Q}{1/2}=2\RD(\GQ)$.
      \emph{Equality case:} Markov is tight when $W_Q$ only takes the values $0$ and $\frac12$. Example: two voters
      with weight $\frac12$ each, which never err at the same time, and such that on a fraction $r$ of the points
      exactly one of them errs. Then $W_Q=\frac12$ on these points, the vote errs on them (tie), $\RD(\BQ)=r$ and
      $\RD(\GQ)=\frac r2$.
\item For fixed $(x,y)$, draw $h,h'\sim Q$ independently. The events ``$h$ errs'' and ``$h'$ errs'' are independent
      with probability $W_Q$ each, so both err with probability $W_Q^2$, and exactly one errs with probability
      $2W_Q(1-W_Q)$; for binary voters, $h$ and $h'$ disagree iff exactly one of them errs. Averaging over $\D$ (Fubini):
      $e_Q=\E W_Q^2$ and $d_Q=\E[2W_Q(1-W_Q)]$. Pointwise $W_Q=W_Q^2+\frac12\cdot2W_Q(1-W_Q)$, hence
      $\RD(\GQ)=e_Q+\frac12d_Q$. Finally Jensen (or $\Var W_Q\geq0$) gives $e_Q=\E W_Q^2\geq(\E W_Q)^2=\RD(\GQ)^2$,
      so $d_Q=2(\RD(\GQ)-e_Q)\leq2\RD(\GQ)(1-\RD(\GQ))$, with equality iff $W_Q$ is almost surely constant.
\item $\P(W_Q\geq\frac12)=\P(W_Q^2\geq\frac14)\leq\frac{\E W_Q^2}{1/4}=4e_Q$ (Markov applied to $W_Q^2$). Using question~2,
      $4e_Q=4\RD(\GQ)-2d_Q$, so $4e_Q\leq2\RD(\GQ)\iff2\RD(\GQ)\leq2d_Q\iff d_Q\geq\RD(\GQ)$: the tandem bound improves on
      the first-order bound exactly when the voters disagree more often than an average voter errs.
\item For $\mu<\frac12$ and $W_Q\geq\frac12$, we have $W_Q-\mu\geq\frac12-\mu>0$, hence $(W_Q-\mu)^2\geq(\frac12-\mu)^2$: the
      event $\{W_Q\geq\frac12\}$ is included in $\{(W_Q-\mu)^2\geq(\frac12-\mu)^2\}$ and Markov gives the bound. Write $R=\E W$,
      $e=\E W^2$, $v=e-R^2$ and $f(\mu)=\frac{e-2\mu R+\mu^2}{(1/2-\mu)^2}$. Then
      \[f'(\mu)=\frac{2(\mu-R)(\frac12-\mu)+2(e-2\mu R+\mu^2)}{(1/2-\mu)^3}=\frac{2\big(e-\frac R2-\mu(\frac12-R)\big)}{(1/2-\mu)^3},\]
      which vanishes at $\mu^\star=\frac{R/2-e}{1/2-R}$ (when $R<\frac12$), is negative before and positive after: $\mu^\star$ is
      the minimum. To compute $f(\mu^\star)$, note that $\mu^\star-R=\frac{R/2-e-R/2+R^2}{1/2-R}=\frac{-v}{1/2-R}$ and
      $\frac12-\mu^\star=\frac{1/4-R/2-R/2+e}{1/2-R}=\frac{v+(1/2-R)^2}{1/2-R}$. The numerator is
      $e-2\mu^\star R+\mu^{\star2}=(\mu^\star-R)^2+v=\frac{v^2}{(1/2-R)^2}+v=\frac{v\,(v+(1/2-R)^2)}{(1/2-R)^2}$, so
      \[f(\mu^\star)=\frac{v\,(v+(1/2-R)^2)}{(1/2-R)^2}\cdot\frac{(1/2-R)^2}{(v+(1/2-R)^2)^2}=\frac{v}{v+(1/2-R)^2},\]
      which is Cantelli's inequality for $W_Q$ with $a=\frac12-R$ (compare with Exercise~\ref{exo:concentration}, question~3).
\item With $M=1-2W_Q$ (the margin), $\Var M=4v$ and $\E M=1-2R$, so $\frac{v}{v+(1/2-R)^2}=\frac{4v}{4v+(1-2R)^2}=\frac{\Var M}{\E M^2}=1-\frac{(\E M)^2}{\E M^2}$.
      Finally $\E M=1-2\RD(\GQ)$ and $\E M^2=1-4\E W+4\E W^2=1-4\RD(\GQ)+4e_Q=1-2d_Q$ (question~2). Hence
      $\RD(\BQ)\leq1-\frac{(1-2\RD(\GQ))^2}{1-2d_Q}$: the C-bound is the best member of the family of question~4.
\item $nW_Q\sim\mathrm{Bin}(n,p)$ since each voter errs independently with probability $p$. Hence $\RD(\GQ)=\E W=p$,
      $e_Q=\E W^2=\Var W+p^2=\frac{p(1-p)}n+p^2$, and $d_Q=2(p-e_Q)=2p(1-p)(1-\frac1n)$.
      For $p=0.3$ and $n=3$, the vote errs iff at least $2$ voters err: $\RD(\BQ)=3p^2(1-p)+p^3=0.189+0.027=0.216$. The
      bounds are: first order $2p=0.6$; $e_Q=0.09+0.07=0.16$, so tandem $=0.64$; $d_Q=2\times0.21\times\frac23=0.28$, so
      C-bound $=1-\frac{(0.4)^2}{1-0.56}=1-\frac{0.16}{0.44}\approx0.636$. For $n=25$: $d_Q=0.42\times0.96=0.4032$,
      $1-2d_Q=0.1936$ and C-bound $=1-\frac{0.16}{0.1936}\approx0.174$, while the true risk is about $0.018$. When $n\to\infty$,
      $1-2d_Q\to1-4p(1-p)=(1-2p)^2$ and the C-bound tends to $0$, like the true risk, at the rate $\frac{4p(1-p)}{n(1-2p)^2}$.
      \emph{Comment:} with $3$ voters, $d_Q=0.28<p=0.3$, and the second-order bounds are worse than the first-order
      one (Proposition 5.7 of the notes); with many independent voters, only the second-order bounds, and
      especially the C-bound, capture the benefit of the vote. The tandem bound tends to $4p^2=0.36$ and never goes
      to $0$.
\end{enumerate}
\end{solution}

%% =====================================================================
\begin{exo}{Linear classifiers with a Gaussian posterior \extype{proofs}}\label{exo:linear}
Voters $h_{\mathbf v}(x)=\sign\langle\mathbf v,x\rangle$, prior $P=\mathcal N(\mathbf 0,I_d)$, posterior $Q_{\mathbf w}=\mathcal N(\mathbf w,I_d)$.
\begin{enumerate}
\item Show that $\BQ(x)=\sign\langle\mathbf w,x\rangle$.
\item Show that $\P_{\mathbf v\sim Q_{\mathbf w}}\big(h_{\mathbf v}(x)\neq y\big)=\Phi\big(-\frac{y\langle\mathbf w,x\rangle}{\|x\|}\big)$ where $\Phi$ is the standard
      Gaussian CDF. Give $\RS(G_{Q_{\mathbf w}})$ and $\KL(Q_{\mathbf w}\|P)$.
\item Show that, for a fixed $C$, minimizing Catoni's bound over $\mathbf w$ amounts to minimizing
      $C\sum_i\Phi(-y_i\langle\mathbf w,x_i\rangle/\|x_i\|)+\frac12\|\mathbf w\|^2$. Compare with the SVM.
\item For $\mathbf w=s\mathbf u$ with $\|\mathbf u\|=1$, describe the evolution of $\RS(G_{Q_{s\mathbf u}})$ and of the KL when $s$ grows.
      Relate to the notion of margin.
\end{enumerate}
\end{exo}

\begin{solution}
\begin{enumerate}
\item Fix $x$. Under $Q_{\mathbf w}$, $\langle\mathbf v,x\rangle$ is Gaussian with mean $\langle\mathbf w,x\rangle$ and variance
      $x^\top I_dx=\|x\|^2$. Hence $\P(\langle\mathbf v,x\rangle>0)=\Phi\big(\frac{\langle\mathbf w,x\rangle}{\|x\|}\big)$ and
      $\E_Q\sign\langle\mathbf v,x\rangle=2\Phi\big(\frac{\langle\mathbf w,x\rangle}{\|x\|}\big)-1$, which is positive iff
      $\langle\mathbf w,x\rangle>0$ (as $\Phi(0)=\frac12$ and $\Phi$ is increasing). Therefore $\BQ(x)=\sign\langle\mathbf w,x\rangle$:
      the vote of infinitely many noisy linear classifiers is the linear classifier at their centre.
\item Similarly $y\langle\mathbf v,x\rangle\sim\mathcal N(y\langle\mathbf w,x\rangle,\|x\|^2)$ (as $y^2=1$), so
      $\P_{\mathbf v}(h_{\mathbf v}(x)\neq y)=\P(y\langle\mathbf v,x\rangle\leq0)=\Phi\big(-\frac{y\langle\mathbf w,x\rangle}{\|x\|}\big)$.
      Averaging over the sample, $\RS(G_{Q_{\mathbf w}})=\frac1m\sum_i\Phi\big(-\frac{y_i\langle\mathbf w,x_i\rangle}{\|x_i\|}\big)$: a
      smooth \emph{probit} loss of the normalized margin $\frac{y_i\langle\mathbf w,x_i\rangle}{\|x_i\|}$. By
      Exercise~\ref{exo:kl}.4 with $\sigma=1$, $\KL(Q_{\mathbf w}\|P)=\frac12\|\mathbf w\|^2$.
\item Catoni's bound reads $\RD(\GQ)\leq\mathcal F_C^{-1}\big(C\RS(\GQ)+\frac{\KL(Q\|P)+\ln(1/\delta)}m\big)$, and
      $\mathcal F_C^{-1}$ is increasing: for fixed $C$ and $\delta$, minimizing the bound over $\mathbf w$ amounts to minimizing
      $C\RS(G_{Q_{\mathbf w}})+\frac{\|\mathbf w\|^2}{2m}$, i.e.\ (multiplying by $m$)
      $C\sum_i\Phi\big(-\frac{y_i\langle\mathbf w,x_i\rangle}{\|x_i\|}\big)+\frac12\|\mathbf w\|^2$. The soft-margin SVM minimizes
      $C\sum_i\max(0,1-y_i\langle\mathbf w,x_i\rangle)+\frac12\|\mathbf w\|^2$: the regularizer is the same (it is the KL to the
      prior), and the hinge loss is replaced by the probit loss. Two differences: the probit loss is bounded (a badly
      misclassified point costs at most $1$, which gives robustness to label noise) but not convex; and it uses the
      normalized margin.
\item For $\mathbf w=s\mathbf u$, the loss of point $i$ is $\Phi(-s\gamma_i)$ with $\gamma_i=\frac{y_i\langle\mathbf u,x_i\rangle}{\|x_i\|}$.
      When $s$ grows, it decreases to $0$ for well-classified points ($\gamma_i>0$), the faster the larger $\gamma_i$, and
      increases to $1$ for misclassified points; at $s=0$ all losses equal $\frac12$. Meanwhile the KL grows as $\frac{s^2}2$.
      If all points have a normalized margin at least $\gamma$, choosing $s\approx\frac{c}{\gamma}$ makes the empirical Gibbs
      risk at most $\Phi(-c)$ at a KL cost of $\frac{c^2}{2\gamma^2}$: large margins allow small KL, which is the
      PAC-Bayesian margin bound (Langford and Shawe-Taylor, 2002). The trade-off in $s$ is the one of Figure~31 of the notes (PBGD).
\end{enumerate}
\end{solution}

%% =====================================================================
\begin{exo}{$(\star)$ The PAC-Bayes-$\lambda$ bound and its minimization \extype{proofs}}\label{exo:lambda}
\begin{enumerate}
\item Prove the refined Pinsker inequality $\kl(q\|p)\geq\frac{(p-q)^2}{2p}$ for $q<p$.
\item Starting from Seeger's bound, show that w.p.\ $\geq1-\delta$, for all $Q$ and all $\lambda\in(0,2)$ simultaneously,
      $\RD(\GQ)\leq\frac{\RS(\GQ)}{1-\lambda/2}+\frac{\KL(Q\|P)+\ln\frac{2\sqrt m}\delta}{\lambda(1-\lambda/2)m}$.
      Why is no union bound over $\lambda$ needed here, contrary to Exercise~\ref{exo:bounds}?
\item For a fixed $Q$, compute the optimal $\lambda$. For a fixed $\lambda$, compute the optimal $Q$ on a finite set of voters.
      Write the alternating minimization algorithm.
\item Explain why the posterior returned by this algorithm on a random forest may give a \emph{worse} majority vote than the uniform one.
\end{enumerate}
\end{exo}

\begin{solution}
\begin{enumerate}
\item Let $0\leq q<p\leq1$ and $g(u)=\kl(u\|p)-\frac{(p-u)^2}{2p}$ for $u\in[0,p]$. Then $g(p)=0$ and
      $g'(u)=\ln\frac{u(1-p)}{p(1-u)}+\frac{p-u}{p}$. Using $\ln t\leq t-1$ with $t=\frac{u(1-p)}{p(1-u)}$:
      $\ln\frac{u(1-p)}{p(1-u)}\leq\frac{u(1-p)-p(1-u)}{p(1-u)}=\frac{u-p}{p(1-u)}\leq\frac{u-p}{p}$, the last step because
      $u-p<0$ and $0<1-u\leq1$. Hence $g'(u)\leq0$: $g$ is non-increasing on $(0,p)$ and $g(q)\geq g(p)=0$ (see also
      Appendix A.3 of the notes).
\item Let $\hat q=\RS(\GQ)$, $p=\RD(\GQ)$ and $\varepsilon=\frac1m[\KL(Q\|P)+\ln\frac{2\sqrt m}\delta]$. On the event of Seeger's bound,
      $\kl(\hat q\|p)\leq\varepsilon$. If $p\leq\hat q$ the claimed bound is trivial. Otherwise the refined Pinsker inequality
      gives $\frac{(p-\hat q)^2}{2p}\leq\varepsilon$, i.e.\ $p-\hat q\leq\sqrt{2p\varepsilon}$. For any $\lambda>0$, the
      arithmetic--geometric inequality $\sqrt{ab}\leq\frac{a+b}2$ with $a=\lambda p$, $b=\frac{2\varepsilon}\lambda$ gives
      $\sqrt{2p\varepsilon}\leq\frac{\lambda p}2+\frac\varepsilon\lambda$. Hence $p(1-\frac\lambda2)\leq\hat q+\frac\varepsilon\lambda$, and for
      $\lambda\in(0,2)$ we can divide by $1-\frac\lambda2>0$. \emph{Why no union bound:} all these inequalities are
      deterministic consequences of the \emph{single} event of Seeger's bound; $\lambda$ appears only in the algebra,
      not in the probabilistic statement. In Exercise~\ref{exo:bounds}, on the contrary, $\lambda$ defined the function
      $\Delta$, hence the event itself.
\item \emph{Optimal $\lambda$ for fixed $Q$.} Let $B(\lambda)=\frac{\hat q}{1-\lambda/2}+\frac{\varepsilon}{\lambda(1-\lambda/2)}$. Setting
      $B'(\lambda)=0$ leads to $\frac{\hat q}{2}\lambda^2=\varepsilon(1-\lambda)$, i.e.\ $\hat q\lambda^2+2\varepsilon\lambda-2\varepsilon=0$, whose
      positive root is $\lambda^\star=\frac{-\varepsilon+\sqrt{\varepsilon^2+2\hat q\varepsilon}}{\hat q}=\frac{2}{\sqrt{2\hat q/\varepsilon+1}+1}$.
      \emph{Optimal $Q$ for fixed $\lambda$.} The bound is an increasing affine function of
      $\RS(\GQ)+\frac{\KL(Q\|P)}{\lambda m}$, minimized by the Gibbs posterior $Q(h)\propto P(h)e^{-\lambda m\RS(h)}$
      (Exercise~\ref{exo:dv}). \emph{Algorithm:} $\lambda\leftarrow1$; repeat $Q\leftarrow$ Gibbs posterior at $\lambda$, then
      $\lambda\leftarrow\lambda^\star(Q)$, until $\lambda$ stabilizes; return $Q$ and the certificate computed with the kl form. Each
      step decreases the bound, which is therefore non-increasing along the iterations (and converges to the global
      minimum under the conditions of Thiemann et al., 2017).
\item The bound only involves the \emph{average} OOB loss $\rho^\top\widehat L$ and the KL. It therefore puts the weight
      on the few trees with the smallest OOB loss, which tend to be similar (they are grown on similar bootstrap
      samples and make similar errors). The effective number of trees collapses (to about $2$ on \texttt{digits}),
      the disagreement vanishes, and the vote loses the benefit of averaging: its risk approaches that of a single
      good tree ($0.150$ instead of $0.044$ on \texttt{digits}, Table 2 of the notes). The bound is excellent for the
      Gibbs classifier, but it is the wrong objective for the vote.
\end{enumerate}
\end{solution}

%% =====================================================================
\begin{exo}{$(\star)$ Out-of-bag samples \extype{proofs}}\label{exo:oob}
Consider bagging with bootstrap samples of size $k$ drawn with replacement among the $m$ training examples.
\begin{enumerate}
\item Compute the probability that a given example is out-of-bag for a given tree; give its limit when $m\to\infty$ for
      $k=m$ and for $k=m/2$. Same question for being out-of-bag for two given trees.
\item Explain why the OOB loss $\widehat L_t$ of tree $t$ is an unbiased estimate of $\RD(h_t)$ given $S_t$, and why a uniform
      prior over the \emph{indices} of the trees is a legitimate PAC-Bayesian prior although the trees depend on $S$.
\item Show that if $\E e^{N\kl(\hat p_N\|p)}\leq2\sqrt N$ for all $N$, then for $n_0\leq N$, $\E e^{n_0\kl(\hat p_N\|p)}\leq2\sqrt{n_0}$.
      Why is this useful?
\end{enumerate}
\end{exo}

\begin{solution}
\begin{enumerate}
\item A bootstrap sample of size $k$ draws $k$ indices uniformly and independently in $\{1,\dots,m\}$. A given example
      is never drawn with probability $(1-\frac1m)^k$. Since $(1-\frac1m)^m\to e^{-1}$, this tends to
      $e^{-1}\approx0.368$ for $k=m$ and to $e^{-1/2}\approx0.607$ for $k=m/2$. For two trees with independent bootstrap
      samples, the example is out-of-bag for both with probability $(1-\frac1m)^{2k}\to e^{-2}\approx0.135$ ($k=m$) or
      $e^{-1}\approx0.368$ ($k=m/2$). Halving the bootstrap size thus almost triples the size of the pairwise OOB sets
      used by the tandem bound, at the price of slightly weaker trees.
\item Condition on the bootstrap sample $S_t$ (and on the internal randomness of the tree). Then $h_t$ is a fixed
      function, and the OOB examples, which were not used to build it, are still i.i.d.\ from $\D$ and independent of
      $h_t$ (they are the examples of $S$ whose index was not drawn; given the drawn indices, their values are
      independent of $S_t$). Hence $\E[\widehat L_t\,|\,S_t]=\RD(h_t)$, and Maurer's lemma applies to $\widehat L_t$
      conditionally on $S_t$. The prior is placed on the \emph{indices} $t\in\{1,\dots,n\}$, not on the trees themselves:
      the uniform distribution on indices is fixed before seeing the data. The dependence of $h_t$ on $S$ only enters the
      third step of the proof, where, for each index $t$, the conditional concentration of $\widehat L_t$ is used.
\item Write $e^{n_0\kl(\hat p_N\|p)}=\big(e^{N\kl(\hat p_N\|p)}\big)^{n_0/N}$. Since $n_0\leq N$, $x\mapsto x^{n_0/N}$ is concave,
      and Jensen gives $\E e^{n_0\kl}\leq\big(\E e^{N\kl}\big)^{n_0/N}\leq(2\sqrt N)^{n_0/N}=\exp\big(n_0\frac{\ln(2\sqrt N)}{N}\big)$.
      The function $N\mapsto\frac{\ln(2\sqrt N)}N$ is decreasing for $N\geq1$ (its derivative is $\frac{1/2-\ln(2\sqrt N)}{N^2}<0$),
      so for $N\geq n_0$ this is at most $\exp\big(n_0\frac{\ln(2\sqrt{n_0})}{n_0}\big)=2\sqrt{n_0}$. \emph{Usefulness:} the OOB sets
      have different sizes; this lemma allows us to use a single sample size $n_0=\min_t|T_t|$ in the complexity term
      while still computing each loss $\widehat L_t$ on \emph{all} the OOB examples of tree $t$.
\end{enumerate}
\end{solution}


%% =====================================================================
\begin{exo}{$(\star)$ Self-bounding learning \extype{comprehension and proofs}}\label{exo:selfbounding}
Let $\mathcal B(Q,S,\delta)$ be a bound such that $\P_S(\forall Q,\ \RD(Q)\leq\mathcal B(Q,S,\delta))\geq1-\delta$.
\begin{enumerate}
\item Let $Q_S$ be the posterior returned by any algorithm run on $S$. Show that $\P_S(\RD(Q_S)\leq\mathcal B(Q_S,S,\delta))\geq1-\delta$.
      Why does this fail for a bound that holds only for a \emph{fixed} predictor (e.g.\ Hoeffding's inequality)?
\item You build $K=5$ random forests with tree depths $1,2,4,8$ and unlimited, and compute for each of them the OOB tandem
      certificate. How must you modify the certificates to be allowed to keep the forest with the smallest one?
      What is the cost in the complexity term?
\item Among the following choices, which are free, which require a union bound, which are forbidden?
      (a) the temperature $\lambda$ of a Gibbs posterior; (b) the parameter $C$ of Catoni's bound; (c) the variance of a
      Gaussian prior chosen on a grid of 10 values; (d) the centre of a Gaussian prior learned on the sample used in the
      bound; (e) the number of gradient iterations when minimizing the bound.
\item Let $p=\klup(q,\varepsilon)$. Show that $\frac{\partial p}{\partial\varepsilon}=\frac{p(1-p)}{p-q}$ and
      $\frac{\partial p}{\partial q}=\frac{p(1-p)}{p-q}\ln\frac{p(1-q)}{q(1-p)}$. Deduce the gradient of the first-order
      certificate $2\,\klup\big(\rho^\top\widehat L,\frac{\KL(\rho\|\pi)+c}{n}\big)$ with respect to $\rho$.
\item Why is the test-set bound usually tighter than a self-bounding certificate for a random forest? Give one
      advantage of the self-bounding approach.
\end{enumerate}
\end{exo}

\begin{solution}
\begin{enumerate}
\item Let $E=\{S:\forall Q,\ \RD(Q)\leq\mathcal B(Q,S,\delta)\}$, with $\P(E)\geq1-\delta$. For $S\in E$ the inequality holds for
      every posterior, in particular for $Q_S$, whatever the way it was computed. Hence
      $\{S:\RD(Q_S)\leq\mathcal B(Q_S,S,\delta)\}\supseteq E$ and has probability $\geq1-\delta$. For a fixed-predictor bound,
      there is one event $E_h$ per predictor, each of probability $\geq1-\delta$, but their intersection can have a much
      smaller probability. The learning algorithm, which looks at $S$, may precisely choose a predictor for which
      $E_h$ fails: among $1000$ random classifiers of risk $0.5$ evaluated on $100$ examples, the one with the smallest
      empirical risk violates the naive Hoeffding bound in $99.8\%$ of the runs (Figure 10 of the notes).
\item Keeping the smallest certificate is a data-dependent choice among $K=5$ bounds. By the union bound
      (Proposition 6.3 of the notes), if each certificate is computed with confidence $\delta/5$, all five hold
      simultaneously with probability $\geq1-\delta$, and so does the smallest one. For the tandem bound, the complexity
      term $\frac{2\KL+\ln(2\sqrt{n_0}/\delta)}{n_0}$ becomes $\frac{2\KL+\ln(2\sqrt{n_0}/\delta)+\ln5}{n_0}$: the price is
      $\ln5\approx1.6$ nats, i.e.\ $0.004$ in the complexity term for $n_0=400$.
\item (a) \emph{Free}: a temperature only indexes a family of posteriors compared with the same prior in the same
      bound. (b) \emph{Union bound}: $C$ defines the function $\Delta$, hence the event; use a grid fixed in advance with
      $\delta/k$. (c) \emph{Union bound}: each variance defines a different prior, hence a different bound; cost $\ln10$.
      (d) \emph{Forbidden}: the prior would depend on the sample used in the bound, and step 3 of the proof fails; it
      becomes legitimate if the centre is learned on a separate part $S_1$ and the bound computed on $S_2$
      (Proposition 6.5). (e) \emph{Free}: any stopping rule produces a posterior, and the bound holds for all posteriors.
\item Let $F(q,p)=\kl(q\|p)-\varepsilon$, so that $F(q,\klup(q,\varepsilon))=0$. We have $F_q=\partial_q\kl(q\|p)=\ln\frac{q(1-p)}{p(1-q)}$
      and $F_p=\frac{p-q}{p(1-p)}\neq0$ (since $p>q$ when $\varepsilon>0$). By the implicit function theorem,
      $\frac{\partial p}{\partial\varepsilon}=-\frac{F_\varepsilon}{F_p}=\frac1{F_p}=\frac{p(1-p)}{p-q}$ and
      $\frac{\partial p}{\partial q}=-\frac{F_q}{F_p}=\frac{p(1-p)}{p-q}\ln\frac{p(1-q)}{q(1-p)}$. For the first-order certificate
      $B(\rho)=2\,\klup(q(\rho),\varepsilon(\rho))$ with $q(\rho)=\rho^\top\widehat L$ and $\varepsilon(\rho)=\frac{\KL(\rho\|\pi)+c}n$,
      we have $\nabla_\rho q=\widehat L$ and $\nabla_\rho\varepsilon=\frac1n\big(\ln\frac\rho\pi+\mathbf1\big)$, so by the chain rule
      \[\nabla_\rho B=2\Big[\frac{\partial p}{\partial q}\,\widehat L+\frac{\partial p}{\partial\varepsilon}\,\frac1n\Big(\ln\frac\rho\pi+\mathbf 1\Big)\Big].\]
      (With a softmax parametrization $\rho=\mathrm{softmax}(\theta)$, multiply by the Jacobian: $\nabla_\theta B=\rho\odot(g-\rho^\top g)$.)
      Numerically, at $q=0.1$, $\varepsilon=0.05$, $p\approx0.220$: $\partial_\varepsilon p\approx1.43$ and $\partial_qp\approx1.33$.
\item The test-set bound certifies the vote itself, with a single kl inversion, no KL term, and no factor $2$ or $4$
      coming from the Gibbs-to-vote relation; on a random forest this relation is the main source of looseness. In
      exchange, the test examples are not used for learning, and a new split is needed for model selection. The
      self-bounding approach uses all the data, provides a learning objective with a guarantee, and allows model
      selection at the price of $\ln K$; with data-dependent priors it can be competitive when enough data are
      available.
\end{enumerate}
\end{solution}

%% =====================================================================
%% =====================================================================
\exoblock{Part III -- Applications and short proofs}

%% =====================================================================
\begin{exo}{A majority vote by hand \extype{application}}\label{exo:votes}
Five voters $h_1,\dots,h_5$ are evaluated on eight labelled points. Their predictions are:
\begin{center}
\renewcommand{\arraystretch}{1.15}
\begin{tabular}{c|cccccccc}
\toprule
point $i$ & 1 & 2 & 3 & 4 & 5 & 6 & 7 & 8\\
\midrule
label $y_i$ & $+$ & $-$ & $+$ & $+$ & $-$ & $+$ & $-$ & $-$\\
\midrule
$h_1(x_i)$ & $-$ & $+$ & $+$ & $+$ & $-$ & $+$ & $-$ & $-$\\
$h_2(x_i)$ & $-$ & $-$ & $-$ & $+$ & $-$ & $+$ & $-$ & $-$\\
$h_3(x_i)$ & $-$ & $-$ & $+$ & $-$ & $-$ & $+$ & $-$ & $-$\\
$h_4(x_i)$ & $+$ & $-$ & $+$ & $+$ & $+$ & $-$ & $-$ & $-$\\
$h_5(x_i)$ & $+$ & $-$ & $+$ & $+$ & $-$ & $+$ & $+$ & $+$\\
\bottomrule
\end{tabular}
\end{center}
We treat the eight points as the whole distribution $\D$ (uniform on the eight points), so that all quantities below are
exact (``oracle'') values. Ties of the vote count as errors.
\begin{enumerate}
\item Give the risk of each voter. For the uniform posterior $Q=(\frac15,\dots,\frac15)$, compute $W_Q(x_i,y_i)$ and the margin
      $M_Q=1-2W_Q$ on each point, then $\RD(\GQ)$ and $\RD(\BQ)$.
\item Compute the joint error $e_Q=\E W_Q^2$ and the disagreement $d_Q=\E[2W_Q(1-W_Q)]$, and check that $\RD(\GQ)=e_Q+\frac12d_Q$.
      Compute the $5\times5$ matrix $D$ of pairwise disagreements $D_{jk}=\P_x(h_j(x)\neq h_k(x))$ and check that
      $d_Q=\sum_{j,k}Q_jQ_kD_{jk}$. Which of these computations need the labels?
\item Compute the first-order bound $2\RD(\GQ)$, the tandem bound $4e_Q$ and the C-bound
      $1-\frac{(1-2\RD(\GQ))^2}{1-2d_Q}$, and compare them with $\RD(\BQ)$. Give the optimal parameter $\mu^\star$ of the
      Chebyshev--Cantelli family (Exercise~\ref{exo:mv}, question~4) and check that it gives the C-bound.
\item Same questions for the posterior $Q'=(0.1,0.1,0.1,0.35,0.35)$. Why do the second-order bounds improve while the
      first-order bound does not move?
\item Same questions for the Dirac posterior $Q''=\delta_{h_1}$. Which bound is the best now? Comment.
\end{enumerate}
\end{exo}

\begin{solution}
\begin{enumerate}
\item Comparing each row with the labels, $h_1$ errs on points $\{1,2\}$, $h_2$ on $\{1,3\}$, $h_3$ on $\{1,4\}$, $h_4$ on $\{5,6\}$
      and $h_5$ on $\{7,8\}$: every voter has risk $\frac28=0.25$. Under the uniform posterior, $W_Q(x_i,y_i)$ is the
      fraction of wrong voters: $W_Q=\frac35=0.6$ on point~1 (three errors) and $\frac15=0.2$ on the seven other points. Hence
      $M_Q=-0.2$ on point~1 and $0.6$ elsewhere. The vote errs only on point~1: $\RD(\BQ)=\frac18=0.125$, and
      $\RD(\GQ)=\E W_Q=\frac{0.6+7\times0.2}8=0.25$, the average of the individual risks as it should.
\item $e_Q=\frac{0.6^2+7\times0.2^2}8=\frac{0.36+0.28}8=0.08$ and $d_Q=\frac{2\times0.6\times0.4+7\times2\times0.2\times0.8}8=\frac{0.48+2.24}8=0.34$;
      indeed $0.08+\frac{0.34}2=0.25$. Two voters disagree exactly on the points where one of them errs and not the
      other. Among $h_1,h_2,h_3$, each pair shares the error on point~1 and disagrees on two points ($D_{12}=D_{13}=D_{23}=0.25$);
      $h_4$ and $h_5$ have errors disjoint from all the others, so any pair involving them disagrees on four points
      ($D_{jk}=0.5$). With $Q_jQ_k=\frac1{25}$ and $20$ ordered pairs $j\neq k$ (six inside $\{h_1,h_2,h_3\}$, fourteen others):
      $\sum_{j,k}Q_jQ_kD_{jk}=\frac{6\times0.25+14\times0.5}{25}=\frac{8.5}{25}=0.34=d_Q$. The matrix $D$, hence $d_Q$, only compares
      rows of predictions: \emph{no label is needed}. $W_Q$, $e_Q$, the risks and the vote risk need the labels.
\item First-order: $2\times0.25=0.5$. Tandem: $4\times0.08=0.32$. C-bound: $1-\frac{(1-0.5)^2}{1-0.68}=1-\frac{0.25}{0.32}=0.219$. All three
      are valid ($\geq0.125$), and the C-bound is the tightest; the tandem bound improves on the first-order one because
      $d_Q=0.34\geq\RD(\GQ)=0.25$. The optimal parameter is $\mu^\star=\frac{\RD(\GQ)/2-e_Q}{1/2-\RD(\GQ)}=\frac{0.125-0.08}{0.25}=0.18$, and
      $\frac{e_Q-2\mu^\star\RD(\GQ)+\mu^{\star2}}{(1/2-\mu^\star)^2}=\frac{0.08-0.09+0.0324}{0.1024}=\frac{0.0224}{0.1024}=0.219$: the C-bound.
\item With $Q'$: $W_{Q'}=0.3$ on point~1, $0.1$ on points~2--4, $0.35$ on points~5--8, so all margins are positive
      ($0.4$, $0.8$, $0.3$) and $\RD(B_{Q'})=0$. The Gibbs risk is unchanged, $\RD(G_{Q'})=\frac{0.3+3\times0.1+4\times0.35}8=0.25$
      (all voters have the same risk), so the first-order bound stays at $0.5$. But
      $e_{Q'}=\frac{0.09+3\times0.01+4\times0.1225}8=\frac{0.61}8=0.07625$ and $d_{Q'}=2(0.25-0.07625)=0.3475$ (check with $D$:
      $6\times0.01\times0.25+12\times0.035\times0.5+2\times0.1225\times0.5=0.015+0.21+0.1225=0.3475$). Tandem: $0.305$; C-bound:
      $1-\frac{0.25}{1-0.695}=0.180$; $\mu^\star=\frac{0.125-0.07625}{0.25}=0.195$. \emph{Why:} $h_1,h_2,h_3$ all err on point~1; moving weight
      towards $h_4,h_5$, whose errors are not shared, decreases the joint error and increases the disagreement. The
      first-order bound only sees the average risk, which is the same for all voters.
\item With $Q''=\delta_{h_1}$, $W\in\{0,1\}$ (the vote is $h_1$): $\RD(B_{Q''})=\RD(G_{Q''})=0.25$, $e=0.25$, $d=0$. First-order: $0.5$;
      tandem: $1$; C-bound: $1-\frac{0.25}{1}=0.75$. The first-order bound is now the best, as predicted by
      Exercise~\ref{exo:quiz}(j): without diversity, second-order bounds are useless.
\end{enumerate}
\emph{What we learn:} the benefit of the vote is invisible to first-order quantities; the joint error and the
disagreement measure it, and the latter can be computed without labels. Even as oracle quantities the bounds remain
pessimistic on such a tiny example (they are worst cases over all distributions of $W_Q$ with given moments).
\end{solution}

%% =====================================================================
\begin{exo}{The Gibbs posterior on four voters \extype{application}}\label{exo:gibbs4}
Four voters have, on a sample of size $m=200$, the empirical risks
$\RS(h_1)=0.10$, $\RS(h_2)=0.12$, $\RS(h_3)=0.15$, $\RS(h_4)=0.30$ (i.e.\ $20$, $24$, $30$ and $60$ errors). The prior $P$
is uniform and $\delta=0.05$. Recall the Gibbs posterior $Q_\lambda(h)\propto P(h)e^{-\lambda m\RS(h)}$ (Exercise~\ref{exo:dv}),
Seeger's bound $\RD(\GQ)\leq\klup\big(\RS(\GQ),\frac{\KL(Q\|P)+\ln(2\sqrt m/\delta)}m\big)$ and McAllester's bound
$\RD(\GQ)\leq\RS(\GQ)+\sqrt{\frac{\KL(Q\|P)+\ln(2\sqrt m/\delta)}{2m}}$. You may use \texttt{kl\_inv\_upper} of \texttt{pbtools}
for the inversions.
\begin{enumerate}
\item Show that $Q_\lambda(h_j)=\frac{e^{-\lambda(k_j-k_1)}}{\sum_le^{-\lambda(k_l-k_1)}}$, where $k_j$ is the number of errors of $h_j$. Why is this
      form numerically preferable? Compute $Q_\lambda$ for $\lambda\in\{0.05,\,0.25,\,1\}$.
\item For $\lambda\in\{0,\,0.05,\,0.25,\,1\}$ and for $\lambda\to+\infty$, compute $\RS(G_{Q_\lambda})$, $\KL(Q_\lambda\|P)$ (use
      $\KL=\ln4-H(Q_\lambda)$, Exercise~\ref{exo:kl}), and both bounds.
\item For $\lambda=0.25$, compute the objective $\RS(\GQ)+\frac{\KL(Q\|P)}{\lambda m}$ for $Q_\lambda$, for the uniform posterior and for
      $\delta_{h_1}$, and check the value $-\frac1{\lambda m}\ln\E_{h\sim P}e^{-\lambda m\RS(h)}$ found in Exercise~\ref{exo:dv}.
\item Which posterior gives the best certificate? Is it legitimate to choose $\lambda$ after computing the bounds?
      Why does the bound favour concentrated posteriors here, and what would this imply for the majority vote?
\end{enumerate}
\end{exo}

\begin{solution}
\begin{enumerate}
\item $Q_\lambda(h_j)=\frac{\frac14e^{-\lambda k_j}}{\sum_l\frac14e^{-\lambda k_l}}$ since $m\RS(h_j)=k_j$; multiplying numerator and denominator by
      $e^{\lambda k_1}$ gives the formula. The largest exponential is now $e^0=1$: no underflow even when $\lambda k_j$ is
      large (with $m=10^5$, $e^{-\lambda k_j}$ would be $0$ in floating point). This is the ``log-sum-exp'' trick used by
      \texttt{softmax} in \texttt{pbtools}. With $k_j-k_1=(0,4,10,40)$:
      \begin{center}
      \begin{tabular}{cccc}\toprule
      $\lambda$ & unnormalized weights & sum & $Q_\lambda$\\\midrule
      $0.05$ & $(1,\ 0.819,\ 0.607,\ 0.135)$ & $2.561$ & $(0.391,\ 0.320,\ 0.237,\ 0.053)$\\
      $0.25$ & $(1,\ 0.368,\ 0.082,\ 4.5\cdot10^{-5})$ & $1.450$ & $(0.690,\ 0.254,\ 0.057,\ 0.000)$\\
      $1$ & $(1,\ 0.018,\ 4.5\cdot10^{-5},\ 4\cdot10^{-18})$ & $1.018$ & $(0.982,\ 0.018,\ 0.000,\ 0.000)$\\\bottomrule
      \end{tabular}
      \end{center}
\item Here $\ln\frac{2\sqrt{200}}{0.05}=\ln565.7=6.338$ and $\varepsilon=\frac{\KL+6.338}{200}$. For $\lambda=0.05$:
      $\RS(G_{Q_\lambda})=0.391\times0.10+0.320\times0.12+0.237\times0.15+0.053\times0.30=0.1288$, $H(Q_\lambda)=1.228$, $\KL=1.386-1.228=0.158$.
      The other columns are obtained in the same way:
      \begin{center}
      \renewcommand{\arraystretch}{1.1}
      \begin{tabular}{lccccc}\toprule
      & $\lambda=0$ (uniform) & $0.05$ & $0.25$ & $1$ & $\infty$ ($\delta_{h_1}$)\\\midrule
      $\RS(G_{Q_\lambda})$ & $0.1675$ & $0.1288$ & $0.1079$ & $0.1004$ & $0.1000$\\
      $\KL(Q_\lambda\|P)$ & $0$ & $0.158$ & $0.619$ & $1.296$ & $1.386$\\
      $\varepsilon$ & $0.0317$ & $0.0325$ & $0.0348$ & $0.0382$ & $0.0386$\\
      Seeger & $0.274$ & $0.229$ & $0.207$ & $0.203$ & $0.203$\\
      McAllester & $0.293$ & $0.256$ & $0.240$ & $0.239$ & $0.239$\\\bottomrule
      \end{tabular}
      \end{center}
      (The Dirac on $h_1$ has $\KL=\ln4$, the Occam price of selecting one voter out of four.) Seeger's bound is always tighter than McAllester's, by $0.02$ to $0.036$, the gap growing as the posterior concentrates on voters with small empirical risk (the risks are far from $\frac12$).
\item With $\lambda m=50$: for $Q_{0.25}$, $0.1079+\frac{0.619}{50}=0.1203$; for the uniform posterior, $0.1675+0=0.1675$; for
      $\delta_{h_1}$, $0.1+\frac{1.386}{50}=0.1277$. The Gibbs posterior is indeed the smallest, and
      $-\frac1{50}\ln\big(\frac14(e^{-5}+e^{-6}+e^{-7.5}+e^{-15})\big)=0.1203$ coincides with its value, as predicted by the
      variational formula (it is a soft minimum of the empirical risks, between $0.10$ and their mean).
\item The bound decreases from $0.274$ (uniform) to about $0.203$ ($\lambda\geq1$); a finer scan gives a flat minimum
      $0.2028$ around $\lambda\approx0.8$, barely better than the Dirac. Choosing $\lambda$ after computing the bounds is
      \emph{free}: $\lambda$ only indexes posteriors, and Seeger's bound holds for all posteriors simultaneously
      (Exercise~\ref{exo:selfbounding}, question~3(a)). With only four voters, the KL can never exceed $\ln4=1.39$ nats,
      small compared with the constant $6.34$, so concentrating on the best voter is cheap and it reduces the
      empirical term by $0.07$. The resulting ``vote'' is essentially $h_1$ alone: all diversity is lost, which is the
      phenomenon of Exercise~\ref{exo:lambda}, question~4. Minimizing a Gibbs-risk bound is the right objective for
      the Gibbs classifier, not for the vote.
\end{enumerate}
\end{solution}

%% =====================================================================
\begin{exo}{Data-dependent priors and sample splitting \extype{short proof and application}}\label{exo:ddprior}
We learn a linear classifier with a Gaussian posterior $Q=\mathcal N(\mathbf w,I_d)$ from a sample $S$ of size $m=1000$,
$\delta=0.05$. To obtain a better prior, we split $S=S_1\cup S_2$ with $|S_1|=m_1$ and $|S_2|=m_2=m-m_1$, learn a vector
$\mathbf w_{S_1}$ on $S_1$ only ($\mathbf w_{S_1}=\mathbf 0$ if $m_1=0$), and use the prior $P_{S_1}=\mathcal N(\mathbf w_{S_1},I_d)$. The
posterior centre $\mathbf w$ is learned on the whole sample $S$.
\begin{enumerate}
\item \emph{(Short proof.)} Show that, with probability at least $1-\delta$ over $(S_1,S_2)$, for all posteriors $Q$,
      $\RD(\GQ)\leq\klup\big(\widehat R_{S_2}(\GQ),\frac{\KL(Q\|P_{S_1})+\ln(2\sqrt{m_2}/\delta)}{m_2}\big)$. Why may $Q$ depend on the
      whole of $S$, while the empirical risk must be computed on $S_2$ only?
\item The following values were observed for four splits (in all cases $\widehat R_{S_2}(\GQ)=0.10$):
      \begin{center}
      \begin{tabular}{ccccc}\toprule
      $m_1$ & $0$ & $200$ & $500$ & $800$\\
      $\|\mathbf w-\mathbf w_{S_1}\|^2$ & $120$ & $30$ & $10$ & $4$\\\bottomrule
      \end{tabular}
      \end{center}
      Compute the KL (Exercise~\ref{exo:kl}) and the certificate for each split.
\item Discuss the trade-off governing the choice of $m_1$.
\item An analyst computes the four certificates and reports the smallest one. What should she report instead?
\item What does the bound become if the posterior is the prior itself, $Q=P_{S_1}$? Relate it to the test-set bound of
      Exercise~\ref{exo:testset}.
\end{enumerate}
\end{exo}

\begin{solution}
\begin{enumerate}
\item Condition on $S_1$. Then $P_{S_1}$ is a fixed distribution, and $S_2\sim\D^{m_2}$ is independent of $S_1$, hence of
      $P_{S_1}$. Seeger's bound applied to the sample $S_2$ with the prior $P_{S_1}$ gives, for every fixed value of $S_1$,
      $\P_{S_2}(E\,|\,S_1)\geq1-\delta$, where $E$ is the event ``for all $Q$, the inequality holds''. Integrating over $S_1$,
      $\P(E)=\E_{S_1}[\P_{S_2}(E\,|\,S_1)]\geq1-\delta$. Since on $E$ the inequality holds \emph{for all} $Q$, it holds for a
      posterior learned on $S_1\cup S_2$. The empirical term, on the contrary, is the variable that concentrates in the
      proof: its exponential moment is computed under $\E_{S_2}\E_{h\sim P_{S_1}}$, which requires $S_2$ to be independent
      of the prior. The empirical risk on $S_1$ is biased in favour of the voters that $P_{S_1}$ puts forward (they were
      chosen because they fit $S_1$), so using $S_1$ in the empirical term would break the proof.
\item By Exercise~\ref{exo:kl}.4 with $\sigma=1$, $\KL=\frac12\|\mathbf w-\mathbf w_{S_1}\|^2$. With $c=\ln(2\sqrt{m_2}/\delta)$:
      \begin{center}
      \begin{tabular}{cccccc}\toprule
      $m_1$ & $m_2$ & $\KL$ & $c$ & $\varepsilon=\frac{\KL+c}{m_2}$ & $\klup(0.10,\varepsilon)$\\\midrule
      $0$ & $1000$ & $60$ & $7.14$ & $0.0671$ & $0.243$\\
      $200$ & $800$ & $15$ & $7.03$ & $0.0275$ & $0.185$\\
      $500$ & $500$ & $5$ & $6.80$ & $0.0236$ & $\mathbf{0.177}$\\
      $800$ & $200$ & $2$ & $6.34$ & $0.0417$ & $0.208$\\\bottomrule
      \end{tabular}
      \end{center}
\item Increasing $m_1$ gives a better prior, closer to the posterior, hence a smaller KL; but it leaves fewer examples
      $m_2$ to estimate the empirical risk and to divide the complexity term. Without splitting, the KL ($60$ nats)
      dominates; with $m_1=800$ the KL is almost gone but $\varepsilon$ is limited by $\frac{c}{m_2}\approx0.03$. The best
      trade-off is intermediate ($m_1=500$ here), as in Figure~15 of the notes. Splitting helps when the KL of the
      data-free prior is large compared with $\ln(2\sqrt m/\delta)$, as for large models.
\item Choosing the split with the data is a choice among $K=4$ bounds (different priors and samples): by the union bound
      (Proposition~6.3 of the notes), each certificate must be computed with $\delta/4$, i.e.\ with $\ln(8\sqrt{m_2}/\delta)$.
      This gives $0.245$, $0.188$, $0.182$ and $0.218$: she can report $0.182$ for $m_1=500$. The price, $\ln4=1.4$ nats, is
      small, but it must be paid (and the list of candidate splits must be fixed before looking at the data).
\item With $Q=P_{S_1}$, $\KL=0$ and the certificate is $\klup\big(\widehat R_{S_2}(G_{P_{S_1}}),\frac{\ln(2\sqrt{m_2}/\delta)}{m_2}\big)$: it is a
      test-set bound on $S_2$ for the (stochastic) classifier $G_{P_{S_1}}$, which was fixed before seeing $S_2$. It is slightly
      looser than the test-set bound of Exercise~\ref{exo:testset} ($\ln(2\sqrt{m_2}/\delta)$ instead of $\ln(1/\delta)$) because it
      also covers all the other posteriors. The two ends of the spectrum are thus pure hold-out ($Q=P_{S_1}$) and pure
      self-bounding ($m_1=0$).
\end{enumerate}
\end{solution}

%% =====================================================================
\begin{exo}{Certifying the balanced error \extype{short proof and application}}\label{exo:balanced}
For a classifier $f$, let $R_+(f)=\P(f(x)\neq y\,|\,y=+1)$ and $R_-(f)=\P(f(x)\neq y\,|\,y=-1)$ be the class-conditional
risks, $\mathrm{BER}(f)=\frac12(R_+(f)+R_-(f))$ the balanced error, and $\pi_+=\P(y=+1)$, so that $\RD(f)=\pi_+R_+(f)+(1-\pi_+)R_-(f)$.
A classifier $f$, learned on another sample, is evaluated on an independent test sample of $n=1000$ examples with
$n_+=100$ positives and $n_-=900$ negatives; it makes $12$ errors on the positives and $18$ on the negatives.
\begin{enumerate}
\item Compute the kl test-set bound on $\RD(f)$. Why is this certificate misleading for this problem?
\item \emph{(Short proof.)} Conditionally on the labels $(y_1,\dots,y_n)$, show that $n_+\widehat R_+(f)\sim\mathrm{Bin}(n_+,R_+(f))$, where
      $\widehat R_+(f)$ is the error rate on the positive test examples. Deduce that with probability at least $1-\delta$,
      simultaneously, $R_+(f)\leq U_+:=\klup\big(\widehat R_+(f),\frac{\ln(2/\delta)}{n_+}\big)$ and
      $R_-(f)\leq U_-:=\klup\big(\widehat R_-(f),\frac{\ln(2/\delta)}{n_-}\big)$; hence
      $\mathrm{BER}(f)\leq\frac{U_++U_-}2$ and $\RD(f)\leq\max(U_+,U_-)$.
\item Compute $U_+$, $U_-$ and the certificate on the balanced error. What does it add to question~1? What could you
      say about $\RD(f)$ if $\pi_+=0.1$ were known?
\item $(\star)$ \emph{(Another route.)} Conditionally on the labels, write the empirical balanced error as a sum
      $\sum_iZ_i$ of independent variables with $Z_i\in[0,c_i]$, $c_i=\frac1{2n_{y_i}}$. Adapting the Chernoff method of
      Exercise~\ref{exo:concentration}, show that $\mathrm{BER}(f)\leq\widehat{\mathrm{BER}}(f)+\sqrt{\frac{\ln(1/\delta)}8\big(\frac1{n_+}+\frac1{n_-}\big)}$
      with probability $\geq1-\delta$. Compute it and compare with question~3.
\item How would you obtain a PAC-Bayesian certificate on the balanced error of the Gibbs classifier $\GQ$, valid for all
      posteriors $Q$? And of the majority vote?
\end{enumerate}
\end{exo}

\begin{solution}
\begin{enumerate}
\item $\widehat R_T(f)=\frac{30}{1000}=0.03$ and $\klup\big(0.03,\frac{\ln20}{1000}\big)=\klup(0.03,0.0030)=0.045$. The certificate is correct
      but dominated by the $900$ negatives: it says nothing about the minority class, on which $f$ errs $12\%$ of the
      time. A classifier predicting $-1$ everywhere would have a test error of $0.10$ and a certificate of $\klup(0.1,0.0030)=0.125$,
      while being useless for the positives.
\item Given the labels, the pairs are independent and $x_i\,|\,y_i\sim\D_{\X|Y=y_i}$. For the $n_+$ indices with $y_i=+1$, the
      indicators $\ind{f(x_i)\neq+1}$ are thus i.i.d.\ Bernoulli with parameter $\P(f(x)\neq y\,|\,y=+1)=R_+(f)$ ($f$ is fixed,
      learned on an independent sample); their sum is $\mathrm{Bin}(n_+,R_+(f))$. The test-set bound (Chernoff) with
      confidence $\delta/2$ gives $\P(R_+(f)>U_+\,|\,y_{1:n})\leq\frac\delta2$, and similarly for the negatives. By the union bound,
      both inequalities hold with conditional probability $\geq1-\delta$, for every value of the labels, hence also
      unconditionally. On this event, $\mathrm{BER}=\frac{R_++R_-}2\leq\frac{U_++U_-}2$, and $\RD(f)=\pi_+R_++(1-\pi_+)R_-$ is a
      convex combination of $R_+$ and $R_-$, hence at most $\max(U_+,U_-)$.
\item $\widehat R_+=0.12$, $\varepsilon_+=\frac{\ln40}{100}=0.0369$, $U_+=\klup(0.12,0.0369)=0.226$; $\widehat R_-=0.02$,
      $\varepsilon_-=\frac{\ln40}{900}=0.0041$, $U_-=\klup(0.02,0.0041)=0.035$. Hence $\mathrm{BER}(f)\leq\frac{0.226+0.035}2=0.131$ (empirical
      value $0.07$). The certificate reveals that the error on the positives may be as large as $23\%$, which the
      global certificate $0.045$ hid; the uncertainty is dominated by the small class ($n_+=100$). If $\pi_+=0.1$ is
      known, $\RD(f)\leq0.1\times0.226+0.9\times0.035=0.054$, slightly worse than $0.045$ (we paid $\delta/2$ and the
      small class twice), but together with a guarantee per class.
\item $\widehat{\mathrm{BER}}=\frac1{2n_+}\sum_{y_i=+1}\ell_i+\frac1{2n_-}\sum_{y_i=-1}\ell_i=\sum_iZ_i$ with $Z_i=\frac{\ell_i}{2n_{y_i}}\in[0,c_i]$, independent given
      the labels, and $\E\sum_iZ_i=\mathrm{BER}$. For $\lambda>0$, by Markov and Hoeffding's lemma applied to each
      $\E Z_i-Z_i$ (range $c_i$),
      \[\P\big(\mathrm{BER}-\widehat{\mathrm{BER}}\geq t\big)\leq e^{-\lambda t}\prod_i\E e^{\lambda(\E Z_i-Z_i)}\leq\exp\Big(-\lambda t+\frac{\lambda^2}8\sum_ic_i^2\Big).\]
      The choice $\lambda=4t/\sum_ic_i^2$ gives $\exp(-2t^2/\sum_ic_i^2)$, and $\sum_ic_i^2=\frac{n_+}{4n_+^2}+\frac{n_-}{4n_-^2}=\frac14\big(\frac1{n_+}+\frac1{n_-}\big)$.
      Setting the probability to $\delta$ gives $t=\sqrt{\frac{\ln(1/\delta)}8\big(\frac1{n_+}+\frac1{n_-}\big)}$. Numerically,
      $t=\sqrt{\frac{2.996}8\times0.0111}=0.0645$ and the bound is $0.07+0.0645=0.135$, slightly worse than $0.131$: it avoids the
      union bound but uses Hoeffding instead of kl, which is costly for the small risk of the negative class.
\item Apply Seeger's bound separately to the positive and negative sub-samples, each with confidence $\delta/2$
      (and the same prior $P$): with probability $\geq1-\delta$, for all $Q$,
      $R_\pm(\GQ)\leq\klup\big(\widehat R_\pm(\GQ),\frac{\KL(Q\|P)+\ln(4\sqrt{n_\pm}/\delta)}{n_\pm}\big)$, and the balanced Gibbs risk is at
      most the average. For the vote, the first-order argument works class by class:
      $R_+(\BQ)=\P(W_Q\geq\frac12\,|\,y=+1)\leq2\E[W_Q\,|\,y=+1]=2R_+(\GQ)$, and similarly for $R_-$, so
      $\mathrm{BER}(\BQ)\leq2\,\mathrm{BER}(\GQ)$.
\end{enumerate}
\emph{What we learn:} a certificate is about one specific risk; for imbalanced problems, certifying each class
separately (union bound) costs little and is much more informative.
\end{solution}

%% =====================================================================
\begin{exo}{Softmax parametrization and the gradient of a certificate \extype{short proofs and application}}\label{exo:softmax}
A posterior $\rho$ on $n_v$ voters is parametrized by logits $\theta\in\R^{n_v}$: $\rho=\mathrm{softmax}(\theta)$, i.e.\
$\rho_t=e^{\theta_t}/\sum_se^{\theta_s}$. Let $\pi$ be the prior, $\widehat L\in[0,1]^{n_v}$ the vector of empirical losses,
$\mathcal B(\rho)$ a differentiable objective and $g=\nabla_\rho\mathcal B$.
\begin{enumerate}
\item Show that $\frac{\partial\rho_t}{\partial\theta_s}=\rho_t(\ind{t=s}-\rho_s)$, i.e.\ the Jacobian is $\mathrm{diag}(\rho)-\rho\rho^\top$.
\item Deduce $\nabla_\theta\mathcal B=\rho\odot\big(g-(\rho^\top g)\mathbf 1\big)$. Show that $\mathbf 1^\top\nabla_\theta\mathcal B=0$ and explain this with the
      invariance $\mathrm{softmax}(\theta+c\mathbf1)=\mathrm{softmax}(\theta)$. Why can $g$ be modified by adding a multiple of $\mathbf1$?
\item Show that $\nabla_\rho\KL(\rho\|\pi)=\ln\frac\rho\pi+\mathbf 1$ and $\nabla_\theta\KL(\rho\|\pi)=\rho\odot\big(\ln\frac\rho\pi-\KL(\rho\|\pi)\mathbf1\big)$. What is
      this gradient at $\rho=\pi$?
\item Let $F(\rho)=\rho^\top\widehat L+\frac{\KL(\rho\|\pi)}{\lambda n}$. Show that $\nabla_\theta F=0$ iff $\rho$ is the Gibbs posterior
      $\rho_t\propto\pi_te^{-\lambda n\widehat L_t}$.
\item Consider the first-order certificate $\mathcal B(\rho)=2\,\klup\big(q(\rho),\varepsilon(\rho)\big)$ with $q(\rho)=\rho^\top\widehat L$ and
      $\varepsilon(\rho)=\frac{\KL(\rho\|\pi)+\ln(2\sqrt n/\delta)}n$. Admitting the derivatives $\frac{\partial p}{\partial q}$ and
      $\frac{\partial p}{\partial\varepsilon}$ of $p=\klup(q,\varepsilon)$ (Exercise~\ref{exo:selfbounding}, question~4), write
      $\nabla_\theta\mathcal B$, and simplify it at $\rho=\pi$.
\item Numerical application: three voters, $\widehat L=(0.1,0.2,0.3)$, $\pi$ uniform, $\theta=\mathbf0$, $n=200$, $\delta=0.05$.
      Compute $\varepsilon$, $p=\klup(q,\varepsilon)$ (with \texttt{kl\_inv\_upper}), $\mathcal B$, the two partial derivatives and
      $\nabla_\theta\mathcal B$. Perform one gradient step $\theta\leftarrow\theta-5\nabla_\theta\mathcal B$ and compute the new certificate.
      Where do repeated steps lead?
\end{enumerate}
\end{exo}

\begin{solution}
\begin{enumerate}
\item Let $Z=\sum_se^{\theta_s}$, so $\partial Z/\partial\theta_s=e^{\theta_s}$. By the quotient rule,
      $\frac{\partial\rho_t}{\partial\theta_s}=\frac{\ind{t=s}e^{\theta_t}Z-e^{\theta_t}e^{\theta_s}}{Z^2}=\ind{t=s}\rho_t-\rho_t\rho_s$. In matrix form
      $J=\mathrm{diag}(\rho)-\rho\rho^\top$, which is symmetric.
\item By the chain rule, $\frac{\partial\mathcal B}{\partial\theta_s}=\sum_tg_t\frac{\partial\rho_t}{\partial\theta_s}=\rho_sg_s-\rho_s\sum_t\rho_tg_t=\rho_s(g_s-\rho^\top g)$,
      which is $\nabla_\theta\mathcal B=J^\top g=\rho\odot(g-(\rho^\top g)\mathbf1)$. Summing over $s$:
      $\mathbf1^\top\nabla_\theta\mathcal B=\rho^\top g-(\rho^\top g)\sum_s\rho_s=0$. Indeed $\theta\mapsto\mathcal B(\mathrm{softmax}(\theta))$ is constant along
      the direction $\mathbf 1$, so its derivative in that direction vanishes. Replacing $g$ by $g+c\mathbf 1$ adds
      $c\rho\odot(\mathbf1-\mathbf1)=0$: only the variations of $\mathcal B$ \emph{within} the simplex matter, and $g$ is defined up
      to the direction $\mathbf1$ (which is why constant terms such as the $+\mathbf1$ of the KL gradient are harmless).
\item $\KL(\rho\|\pi)=\sum_t\rho_t\ln\rho_t-\rho_t\ln\pi_t$, so $\frac{\partial\KL}{\partial\rho_t}=\ln\rho_t+1-\ln\pi_t$. Then
      $\rho^\top g=\KL(\rho\|\pi)+1$, and question~2 gives $\nabla_\theta\KL=\rho\odot\big(\ln\frac\rho\pi+\mathbf1-(\KL+1)\mathbf1\big)=\rho\odot\big(\ln\frac\rho\pi-\KL\,\mathbf1\big)$.
      At $\rho=\pi$, $\ln\frac\rho\pi=\mathbf0$ and $\KL=0$: the gradient vanishes. The KL is minimal at the prior, so moving
      slightly away from $\pi$ costs nothing at first order.
\item Here $g=\widehat L+\frac1{\lambda n}\big(\ln\frac\rho\pi+\mathbf1\big)$, so by questions~2--3,
      $\nabla_\theta F=\rho\odot\big(\widehat L-(\rho^\top\widehat L)\mathbf1+\frac1{\lambda n}(\ln\frac\rho\pi-\KL\,\mathbf1)\big)$. Since all $\rho_t>0$, this vanishes iff
      $\widehat L_t+\frac1{\lambda n}\ln\frac{\rho_t}{\pi_t}$ does not depend on $t$, i.e.\ $\ln\frac{\rho_t}{\pi_t}=-\lambda n\widehat L_t+\text{const}$, i.e.\
      $\rho_t\propto\pi_te^{-\lambda n\widehat L_t}$. We recover the Gibbs posterior of Exercise~\ref{exo:dv} without Lagrange
      multiplier: the softmax parametrization handles the constraint automatically.
\item With $p=\klup(q,\varepsilon)$, the chain rule gives $g=\nabla_\rho\mathcal B=2\big[\frac{\partial p}{\partial q}\widehat L+\frac{\partial p}{\partial\varepsilon}\frac1n(\ln\frac\rho\pi+\mathbf 1)\big]$, and by
      questions~2--3,
      \[\nabla_\theta\mathcal B=2\Big[\frac{\partial p}{\partial q}\,\rho\odot\big(\widehat L-q\mathbf 1\big)+\frac{\partial p}{\partial\varepsilon}\,\frac1n\,\rho\odot\Big(\ln\frac\rho\pi-\KL(\rho\|\pi)\mathbf1\Big)\Big].\]
      At $\rho=\pi$ the second term vanishes and $\nabla_\theta\mathcal B=2\frac{\partial p}{\partial q}\,\pi\odot(\widehat L-q\mathbf 1)$: the logits of
      the voters better than average increase, the others decrease.
\item $q=0.2$, $\KL=0$, $\ln\frac{2\sqrt{200}}{0.05}=6.338$, $\varepsilon=0.0317$, $p=\klup(0.2,0.0317)=0.312$ and $\mathcal B=0.624$. Then
      $\frac{\partial p}{\partial\varepsilon}=\frac{p(1-p)}{p-q}=\frac{0.312\times0.688}{0.112}=1.92$ and
      $\frac{\partial p}{\partial q}=1.92\times\ln\frac{p(1-q)}{q(1-p)}=1.92\times\ln1.813=1.92\times0.595=1.14$. Hence
      $\nabla_\theta\mathcal B=2\times1.14\times\frac13(-0.1,0,0.1)=(-0.076,0,0.076)$ (this agrees with
      \texttt{certificate\_and\_grad} of \texttt{pbtools}). After the step, $\theta=(0.380,0,-0.380)$, $\rho=(0.465,0.318,0.217)$,
      $q=0.175$, $\KL=\ln3-H(\rho)=0.047$, $\varepsilon=0.0319$ and $\mathcal B=2\klup(0.175,0.0319)=0.568$: the certificate decreased by $0.056$,
      close to the first-order prediction $5\|\nabla_\theta\mathcal B\|^2=0.058$. Repeated steps keep moving the weight towards
      $h_1$, until the growth of the KL balances the decrease of $q$: the first-order certificate favours concentrated
      posteriors (Exercises~\ref{exo:lambda} and~\ref{exo:gibbs4}). For the tandem certificate, the gradient contains
      $2\widehat{\mathbf L}\rho$ instead of $\widehat L$, which penalizes voters that err together and preserves diversity.
\end{enumerate}
\end{solution}

\end{document}
